
\documentclass{article}
\usepackage{colortbl}
\usepackage{makecell}
\usepackage{multirow}
\usepackage{supertabular}

\begin{document}

\newcounter{utterance}

\centering \large Interaction Transcript for game `clean\_up', experiment `1\_medium\_7obj\_tr', episode 0 with gemma{-}4{-}26b{-}a4b{-}it.
\vspace{24pt}

{ \footnotesize  \setcounter{utterance}{1}
\setlength{\tabcolsep}{0pt}
\begin{supertabular}{c@{$\;$}|p{.15\linewidth}@{}p{.15\linewidth}p{.15\linewidth}p{.15\linewidth}p{.15\linewidth}p{.15\linewidth}}
    \# & $\;$A & \multicolumn{4}{c}{Game Master} & $\;\:$B\\
    \hline

    \theutterance \stepcounter{utterance}  
    & & \multicolumn{4}{p{0.6\linewidth}}{
        \cellcolor[rgb]{0.9,0.9,0.9}{
            \makecell[{{p{\linewidth}}}]{
                \texttt{\tiny{[P1$\langle$GM]}}
                \texttt{Ben senin oyun yöneticinim ve aşağıdaki ızgarayı oyun tahtası olarak kullandığınız işbirlikçi bir oyun oynuyorsunuz.} \\
\texttt{```} \\
\texttt{ 1234567} \\
\texttt{╔════╤══╗} \\
\texttt{║◌◌◌O│PM║ 1} \\
\texttt{║◌◌T◌├──╢ 2} \\
\texttt{║◌◌◌◌│◌◌║ 3} \\
\texttt{║S┌──┼──╢ 4} \\
\texttt{║◌│◌A│◌◌║ 5} \\
\texttt{╟─┴┬─┤◌◌║ 6} \\
\texttt{║◌◌│◌│D◌║ 7} \\
\texttt{╚══╧═╧══╝} \\
\\ 
\texttt{```} \\
\\ 
\texttt{* Üst kenarda x koordinatları sağa doğru artar, sağ kenarda ise y koordinatları aşağı doğru artar.} \\
\texttt{* Aşağıdaki nesneler oyun tahtasına rastgele yerleştirilmiştir: 'P', 'O', 'T', 'S', 'D', 'A', 'M'.} \\
\\ 
\\ 
\texttt{Diğer oyuncu, nesnelerin farklı rastgele konumlara yerleştirildiği oyun tahtasının bir varyasyonunu görür. Sen diğer oyuncunun tahtasını göremezsin ve onlar da seninkini göremez.} \\
\texttt{**Amaç:**} \\
\\ 
\texttt{Her iki oyuncunun da, kendi oyun tahtalarındaki nesneleri hareket ettirerek aynı türden nesnelerin aynı koordinatlara gelmesini sağlaması gerekir. Ortak bir hedef düzen üzerinde anlaşmak için diğer oyuncuyla iletişim kurmalısınız.} \\
\\ 
\\ 
\texttt{**Kurallar:**} \\
\\ 
\texttt{* Her turda, aşağıdaki iki komuttan yalnızca birini gönderebilirsin:} \\
\texttt{ 1. `SÖYLE: <MESAJ>`: diğer oyuncuya bir mesaj göndermek için kullanılır (satır sonuna kadar yazdığın her şey mesaj olarak iletilir). Ben bu mesajı partnerine ileteceğim.} \\
\texttt{2. `HAMLE: <NESNE>, (<X>, <Y>)`: bir nesneyi yeni bir konuma taşımak için kullanılır. Burada <X> sütunu, <Y> ise satırı ifade eder. Yaptığın hamlenin başarılı olup olmadığını sana söyleyeceğim.} \\
\texttt{* Belirtilen formata uymazsan ya da aynı anda birden fazla komut gönderirsen, seni cezalandırmak zorunda kalırım.} \\
\texttt{* Her iki oyuncu da 16 sayısından fazla ceza alırsa, oyunu birlikte kaybedersiniz.} \\
\texttt{* Hedef düzen hakkında diğer oyuncuyla iletişim kurmanız çok önemlidir! Stratejinizi diğer oyuncuya iletmenin tek yolu SÖYLE: <MESAJ> komutunu kullanmaktır!} \\
\\ 
\\ 
\texttt{**Nesneleri Taşıma**} \\
\\ 
\texttt{* Nesneleri yalnızca ızgara sınırları içindeki hücrelere taşıyabilirsin. Hedef hücre boş olmalıdır; yani yalnızca ‘◌’ sembolünü içermelidir.} \\
\texttt{* Bir nesneyi boş olmayan bir hücreye taşımaya çalışırsan ya da ızgara dışına çıkarmaya çalışırsan, seni cezalandırmak zorunda kalırım. Ardından tekrar deneme hakkın olur.} \\
\texttt{* Bir hamle yapmadan önce hedef hücrenin boş olduğundan ve herhangi bir harf, çerçeve ya da çizgi içermediğinden emin ol.} \\
\\ 
\\ 
\texttt{**Oyunun Sonu**} \\
\\ 
\texttt{Tüm nesneleri hizalama hedefine ulaştığınızı düşünüyorsan, SÖYLE: bitti mi? göndererek diğer oyuncudan oyunu bitirmesini isteyebilirsin. Eğer diğer oyuncu senden oyunu bitirmeni isterse ve sen SÖYLE: bitti! ile yanıt verirsen, oyun sona erer.} \\
\\ 
\texttt{Oyun, 28 tur içinde sona ererse her iki oyuncu da kazanır. Bir tur, iki oyuncunun da geçerli bir komut göndermesi olarak tanımlanır.} \\
\\ 
\texttt{**Puanlama:**} \\
\\ 
\texttt{Aynı nesneler her iki oyun tahtasında ne kadar yakın konumlanırsa, o kadar fazla puan kazanırsınız. Cezalar puanınızı düşürür. Rekoru kırabilir misiniz?} \\
\\ 
\\ 
\texttt{Oyuna başlamak için diğer oyuncuya bir mesaj gönder!} \\
            }
        }
    }
    & \\ \\

    \theutterance \stepcounter{utterance}  
    & \multicolumn{4}{p{0.6\linewidth}}{
        \cellcolor[rgb]{0.9,0.9,0.9}{
            \makecell[{{p{\linewidth}}}]{
                \texttt{\tiny{[P1$\rangle$GM]}}
                \texttt{SÖYLE: Merhaba! Seninle işbirliği yapmaya hazırım. Ortak bir hedef düzen belirlemek için önce elimizdeki nesnelerin konumlarını paylaşalım mı? Ben kendi tahtamdaki nesnelerin koordinatlarını aşağıya yazıyorum. Sen de kendi tahtandakileri yazarsan, ortak bir hedef düzen üzerinde hızlıca anlaşabiliriz.} \\
\\ 
\texttt{Benim tahtamdaki nesneler:} \\
\texttt{O: (4, 1)} \\
\texttt{P: (5, 1)} \\
\texttt{M: (6, 1)} \\
\texttt{T: (3, 2)} \\
\texttt{S: (1, 4)} \\
\texttt{A: (4, 5)} \\
\texttt{D: (5, 7)} \\
            }
        }
    }
    & & \\ \\

    \theutterance \stepcounter{utterance}  
    & & & \multicolumn{2}{p{0.3\linewidth}}{
        \cellcolor[rgb]{0.9,0.9,0.9}{
            \makecell[{{p{\linewidth}}}]{
                \texttt{\tiny{[GM$|$GM]}}
                \texttt{SÖYLE: Merhaba! Seninle işbirliği yapmaya hazırım. Ortak bir hedef düzen belirlemek için önce elimizdeki nesnelerin konumlarını paylaşalım mı? Ben kendi tahtamdaki nesnelerin koordinatlarını aşağıya yazıyorum. Sen de kendi tahtandakileri yazarsan, ortak bir hedef düzen üzerinde hızlıca anlaşabiliriz.} \\
\\ 
\texttt{Benim tahtamdaki nesneler:} \\
\texttt{O: (4, 1)} \\
\texttt{P: (5, 1)} \\
\texttt{M: (6, 1)} \\
\texttt{T: (3, 2)} \\
\texttt{S: (1, 4)} \\
\texttt{A: (4, 5)} \\
\texttt{D: (5, 7)} \\
            }
        }
    }
    & & \\ \\

    \theutterance \stepcounter{utterance}  
    & & & \multicolumn{2}{p{0.3\linewidth}}{
        \cellcolor[rgb]{0.9,0.9,0.9}{
            \makecell[{{p{\linewidth}}}]{
                \texttt{\tiny{[GM$|$GM]}}
                \texttt{Invalid format: tail is not empty: '} \\
\\ 
\texttt{Benim tahtamdaki nesneler:} \\
\texttt{O: (4, 1)} \\
\texttt{P: (5, 1)} \\
\texttt{M: (6, 1)} \\
\texttt{T: (3, 2)} \\
\texttt{S: (1, 4)} \\
\texttt{A: (4, 5)} \\
\texttt{D: (5, 7)} \\
            }
        }
    }
    & & \\ \\

    \theutterance \stepcounter{utterance}  
    & & \multicolumn{4}{p{0.6\linewidth}}{
        \cellcolor[rgb]{0.9,0.9,0.9}{
            \makecell[{{p{\linewidth}}}]{
                \texttt{\tiny{[P1$\langle$GM]}}
                \texttt{Ceza: Mesajında komuttan sonra hiçbir şey bulunmamalı!} \\
\\ 
\texttt{Yanıtının yalnızca SÖYLE: <MESAJ> ya da HAMLE: <NESNE>, (<X>, <Y>) komutlarından birini içerdiğinden emin ol; bunun dışında hiçbir şey yazma!} \\
\\ 
\texttt{Toplamda 16 cezanın 1 kadarını biriktirdiniz.Lütfen tekrar deneyin!} \\
            }
        }
    }
    & \\ \\

    \theutterance \stepcounter{utterance}  
    & \multicolumn{4}{p{0.6\linewidth}}{
        \cellcolor[rgb]{0.9,0.9,0.9}{
            \makecell[{{p{\linewidth}}}]{
                \texttt{\tiny{[P1$\rangle$GM]}}
                \texttt{SÖYLE: Merhaba! Ortak bir hedef düzen belirlemek için nesnelerin konumlarını paylaşalım. Benim tahtamdaki nesneler: O (4,1), P (5,1), M (6,1), T (3,2), S (1,4), A (4,5), D (5,7). Seninkileri de bekliyorum.} \\
            }
        }
    }
    & & \\ \\

    \theutterance \stepcounter{utterance}  
    & & & \multicolumn{2}{p{0.3\linewidth}}{
        \cellcolor[rgb]{0.9,0.9,0.9}{
            \makecell[{{p{\linewidth}}}]{
                \texttt{\tiny{[GM$|$GM]}}
                \texttt{SÖYLE: Merhaba! Ortak bir hedef düzen belirlemek için nesnelerin konumlarını paylaşalım. Benim tahtamdaki nesneler: O (4,1), P (5,1), M (6,1), T (3,2), S (1,4), A (4,5), D (5,7). Seninkileri de bekliyorum.} \\
            }
        }
    }
    & & \\ \\

    \theutterance \stepcounter{utterance}  
    & & \multicolumn{4}{p{0.6\linewidth}}{
        \cellcolor[rgb]{0.9,0.9,0.9}{
            \makecell[{{p{\linewidth}}}]{
                \texttt{\tiny{[GM$\rangle$P2]}}
                \texttt{Ben senin oyun yöneticinim ve aşağıdaki ızgarayı oyun tahtası olarak kullandığınız işbirlikçi bir oyun oynuyorsunuz.} \\
\texttt{```} \\
\texttt{ 1234567} \\
\texttt{╔════╤══╗} \\
\texttt{║◌◌◌◌│O◌║ 1} \\
\texttt{║P◌M◌├──╢ 2} \\
\texttt{║◌S◌◌│T◌║ 3} \\
\texttt{║◌┌──┼──╢ 4} \\
\texttt{║A│◌◌│◌◌║ 5} \\
\texttt{╟─┴┬─┤◌D║ 6} \\
\texttt{║◌◌│◌│◌◌║ 7} \\
\texttt{╚══╧═╧══╝} \\
\\ 
\texttt{```} \\
\\ 
\texttt{* Üst kenarda x koordinatları sağa doğru artar, sağ kenarda ise y koordinatları aşağı doğru artar.} \\
\texttt{* Aşağıdaki nesneler oyun tahtasına rastgele yerleştirilmiştir: 'P', 'O', 'T', 'S', 'D', 'A', 'M'.} \\
\\ 
\\ 
\texttt{Diğer oyuncu, nesnelerin farklı rastgele konumlara yerleştirildiği oyun tahtasının bir varyasyonunu görür. Sen diğer oyuncunun tahtasını göremezsin ve onlar da seninkini göremez.} \\
\texttt{**Amaç:**} \\
\\ 
\texttt{Her iki oyuncunun da, kendi oyun tahtalarındaki nesneleri hareket ettirerek aynı türden nesnelerin aynı koordinatlara gelmesini sağlaması gerekir. Ortak bir hedef düzen üzerinde anlaşmak için diğer oyuncuyla iletişim kurmalısınız.} \\
\\ 
\\ 
\texttt{**Kurallar:**} \\
\\ 
\texttt{* Her turda, aşağıdaki iki komuttan yalnızca birini gönderebilirsin:} \\
\texttt{ 1. `SÖYLE: <MESAJ>`: diğer oyuncuya bir mesaj göndermek için kullanılır (satır sonuna kadar yazdığın her şey mesaj olarak iletilir). Ben bu mesajı partnerine ileteceğim.} \\
\texttt{2. `HAMLE: <NESNE>, (<X>, <Y>)`: bir nesneyi yeni bir konuma taşımak için kullanılır. Burada <X> sütunu, <Y> ise satırı ifade eder. Yaptığın hamlenin başarılı olup olmadığını sana söyleyeceğim.} \\
\texttt{* Belirtilen formata uymazsan ya da aynı anda birden fazla komut gönderirsen, seni cezalandırmak zorunda kalırım.} \\
\texttt{* Her iki oyuncu da 16 sayısından fazla ceza alırsa, oyunu birlikte kaybedersiniz.} \\
\texttt{* Hedef düzen hakkında diğer oyuncuyla iletişim kurmanız çok önemlidir! Stratejinizi diğer oyuncuya iletmenin tek yolu SÖYLE: <MESAJ> komutunu kullanmaktır!} \\
\\ 
\\ 
\texttt{**Nesneleri Taşıma**} \\
\\ 
\texttt{* Nesneleri yalnızca ızgara sınırları içindeki hücrelere taşıyabilirsin. Hedef hücre boş olmalıdır; yani yalnızca ‘◌’ sembolünü içermelidir.} \\
\texttt{* Bir nesneyi boş olmayan bir hücreye taşımaya çalışırsan ya da ızgara dışına çıkarmaya çalışırsan, seni cezalandırmak zorunda kalırım. Ardından tekrar deneme hakkın olur.} \\
\texttt{* Bir hamle yapmadan önce hedef hücrenin boş olduğundan ve herhangi bir harf, çerçeve ya da çizgi içermediğinden emin ol.} \\
\\ 
\\ 
\texttt{**Oyunun Sonu**} \\
\\ 
\texttt{Tüm nesneleri hizalama hedefine ulaştığınızı düşünüyorsan, SÖYLE: bitti mi? göndererek diğer oyuncudan oyunu bitirmesini isteyebilirsin. Eğer diğer oyuncu senden oyunu bitirmeni isterse ve sen SÖYLE: bitti! ile yanıt verirsen, oyun sona erer.} \\
\\ 
\texttt{Oyun, 28 tur içinde sona ererse her iki oyuncu da kazanır. Bir tur, iki oyuncunun da geçerli bir komut göndermesi olarak tanımlanır.} \\
\\ 
\texttt{**Puanlama:**} \\
\\ 
\texttt{Aynı nesneler her iki oyun tahtasında ne kadar yakın konumlanırsa, o kadar fazla puan kazanırsınız. Cezalar puanınızı düşürür. Rekoru kırabilir misiniz?} \\
\\ 
\texttt{Diğer oyuncu şu mesajı göndererek oyunu başlattı:} \\
\\ 
\texttt{"Merhaba! Ortak bir hedef düzen belirlemek için nesnelerin konumlarını paylaşalım. Benim tahtamdaki nesneler: O (4,1), P (5,1), M (6,1), T (3,2), S (1,4), A (4,5), D (5,7). Seninkileri de bekliyorum."} \\
\\ 
\texttt{İlk komutun ne?} \\
            }
        }
    }
    & \\ \\

    \theutterance \stepcounter{utterance}  
    & & & \multicolumn{4}{p{0.6\linewidth}}{
        \cellcolor[rgb]{0.9,0.9,0.9}{
            \makecell[{{p{\linewidth}}}]{
                \texttt{\tiny{[GM$\langle$P2]}}
                \texttt{SÖYLE: Merhaba! Konumları paylaştığın için teşekkürler. Benim tahtamdaki nesnelerin konumları şöyle: P (1,2), O (5,1), M (3,2), S (2,3), T (5,3), A (1,5), D (6,6). Ortak bir hedef düzen belirlemek için önce nesneleri bir araya toplamayı mı yoksa belirli bir geometrik şekil oluşturmayı mı önerirsin?} \\
            }
        }
    }
     \\ \\

    \theutterance \stepcounter{utterance}  
    & & & \multicolumn{2}{p{0.3\linewidth}}{
        \cellcolor[rgb]{0.9,0.9,0.9}{
            \makecell[{{p{\linewidth}}}]{
                \texttt{\tiny{[GM$|$GM]}}
                \texttt{SÖYLE: Merhaba! Konumları paylaştığın için teşekkürler. Benim tahtamdaki nesnelerin konumları şöyle: P (1,2), O (5,1), M (3,2), S (2,3), T (5,3), A (1,5), D (6,6). Ortak bir hedef düzen belirlemek için önce nesneleri bir araya toplamayı mı yoksa belirli bir geometrik şekil oluşturmayı mı önerirsin?} \\
            }
        }
    }
    & & \\ \\

    \theutterance \stepcounter{utterance}  
    & & \multicolumn{4}{p{0.6\linewidth}}{
        \cellcolor[rgb]{0.9,0.9,0.9}{
            \makecell[{{p{\linewidth}}}]{
                \texttt{\tiny{[P1$\langle$GM]}}
                \texttt{Mesajın diğer oyuncuya iletildi.} \\
\\ 
\\ 
\texttt{Şu anda en fazla 28 turdan 1. turu oynuyorsunuz.} \\
\texttt{Toplamda 16 cezanın 1 kadarını biriktirdiniz.Diğer oyuncu şu mesajı gönderdi:} \\
\\ 
\\ 
\texttt{"Merhaba! Konumları paylaştığın için teşekkürler. Benim tahtamdaki nesnelerin konumları şöyle: P (1,2), O (5,1), M (3,2), S (2,3), T (5,3), A (1,5), D (6,6). Ortak bir hedef düzen belirlemek için önce nesneleri bir araya toplamayı mı yoksa belirli bir geometrik şekil oluşturmayı mı önerirsin?"} \\
\\ 
\\ 
\texttt{Bir sonraki komutun nedir?} \\
            }
        }
    }
    & \\ \\

    \theutterance \stepcounter{utterance}  
    & \multicolumn{4}{p{0.6\linewidth}}{
        \cellcolor[rgb]{0.9,0.9,0.9}{
            \makecell[{{p{\linewidth}}}]{
                \texttt{\tiny{[P1$\rangle$GM]}}
                \texttt{SÖYLE: Teşekkürler! Nesneleri belirli bir geometrik şekil yerine, koordinatları daha düzenli ve toplu bir hale getirmeye ne dersin? Örneğin, nesneleri ızgaranın üst kısmında veya merkezinde bir blok oluşturacak şekilde hizalayabiliriz. Senin için uygunsa, önce nesneleri (1,1) civarında bir küme haline getirmeyi deneyelim mi?} \\
            }
        }
    }
    & & \\ \\

    \theutterance \stepcounter{utterance}  
    & & & \multicolumn{2}{p{0.3\linewidth}}{
        \cellcolor[rgb]{0.9,0.9,0.9}{
            \makecell[{{p{\linewidth}}}]{
                \texttt{\tiny{[GM$|$GM]}}
                \texttt{SÖYLE: Teşekkürler! Nesneleri belirli bir geometrik şekil yerine, koordinatları daha düzenli ve toplu bir hale getirmeye ne dersin? Örneğin, nesneleri ızgaranın üst kısmında veya merkezinde bir blok oluşturacak şekilde hizalayabiliriz. Senin için uygunsa, önce nesneleri (1,1) civarında bir küme haline getirmeyi deneyelim mi?} \\
            }
        }
    }
    & & \\ \\

    \theutterance \stepcounter{utterance}  
    & & \multicolumn{4}{p{0.6\linewidth}}{
        \cellcolor[rgb]{0.9,0.9,0.9}{
            \makecell[{{p{\linewidth}}}]{
                \texttt{\tiny{[GM$\rangle$P2]}}
                \texttt{Mesajın diğer oyuncuya iletildi.} \\
\\ 
\\ 
\texttt{Şu anda en fazla 28 turdan 2. turu oynuyorsunuz.} \\
\texttt{Toplamda 16 cezanın 1 kadarını biriktirdiniz.Diğer oyuncu şu mesajı gönderdi:} \\
\\ 
\\ 
\texttt{"Teşekkürler! Nesneleri belirli bir geometrik şekil yerine, koordinatları daha düzenli ve toplu bir hale getirmeye ne dersin? Örneğin, nesneleri ızgaranın üst kısmında veya merkezinde bir blok oluşturacak şekilde hizalayabiliriz. Senin için uygunsa, önce nesneleri (1,1) civarında bir küme haline getirmeyi deneyelim mi?"} \\
\\ 
\\ 
\texttt{Bir sonraki komutun nedir?} \\
            }
        }
    }
    & \\ \\

    \theutterance \stepcounter{utterance}  
    & & & \multicolumn{4}{p{0.6\linewidth}}{
        \cellcolor[rgb]{0.9,0.9,0.9}{
            \makecell[{{p{\linewidth}}}]{
                \texttt{\tiny{[GM$\langle$P2]}}
                \texttt{SÖYLE: Harika bir fikir! Nesneleri (1,1) civarında toplu bir küme haline getirmek koordinasyon açısından çok mantıklı. Ben de bu öneriye katılıyorum. İlk hamle olarak, benim tahtamdaki en uzak nesnelerden birini bu bölgeye yaklaştırmaya başlayabilirim. Sen de kendi tahtandaki nesneleri bu merkeze doğru hareket ettirmeye başlayabilirsin. İlk olarak hangi nesneyi hareket ettirmek istersin?} \\
            }
        }
    }
     \\ \\

    \theutterance \stepcounter{utterance}  
    & & & \multicolumn{2}{p{0.3\linewidth}}{
        \cellcolor[rgb]{0.9,0.9,0.9}{
            \makecell[{{p{\linewidth}}}]{
                \texttt{\tiny{[GM$|$GM]}}
                \texttt{SÖYLE: Harika bir fikir! Nesneleri (1,1) civarında toplu bir küme haline getirmek koordinasyon açısından çok mantıklı. Ben de bu öneriye katılıyorum. İlk hamle olarak, benim tahtamdaki en uzak nesnelerden birini bu bölgeye yaklaştırmaya başlayabilirim. Sen de kendi tahtandaki nesneleri bu merkeze doğru hareket ettirmeye başlayabilirsin. İlk olarak hangi nesneyi hareket ettirmek istersin?} \\
            }
        }
    }
    & & \\ \\

    \theutterance \stepcounter{utterance}  
    & & \multicolumn{4}{p{0.6\linewidth}}{
        \cellcolor[rgb]{0.9,0.9,0.9}{
            \makecell[{{p{\linewidth}}}]{
                \texttt{\tiny{[P1$\langle$GM]}}
                \texttt{Mesajın diğer oyuncuya iletildi.} \\
\\ 
\\ 
\texttt{Şu anda en fazla 28 turdan 2. turu oynuyorsunuz.} \\
\texttt{Toplamda 16 cezanın 1 kadarını biriktirdiniz.Diğer oyuncu şu mesajı gönderdi:} \\
\\ 
\\ 
\texttt{"Harika bir fikir! Nesneleri (1,1) civarında toplu bir küme haline getirmek koordinasyon açısından çok mantıklı. Ben de bu öneriye katılıyorum. İlk hamle olarak, benim tahtamdaki en uzak nesnelerden birini bu bölgeye yaklaştırmaya başlayabilirim. Sen de kendi tahtandaki nesneleri bu merkeze doğru hareket ettirmeye başlayabilirsin. İlk olarak hangi nesneyi hareket ettirmek istersin?"} \\
\\ 
\\ 
\texttt{Bir sonraki komutun nedir?} \\
            }
        }
    }
    & \\ \\

    \theutterance \stepcounter{utterance}  
    & \multicolumn{4}{p{0.6\linewidth}}{
        \cellcolor[rgb]{0.9,0.9,0.9}{
            \makecell[{{p{\linewidth}}}]{
                \texttt{\tiny{[P1$\rangle$GM]}}
                \texttt{SÖYLE: Harika! O zaman ben ilk olarak en uzak nesnelerden biri olan D (5,7) nesnesini merkeze yaklaştırmak için hareket ettireceğim. Ben D nesnesini (4,4) konumuna taşımayı deniyorum. Sen de kendi tahtanda uygun bir nesneyi merkeze doğru hareket ettirebilirsin.} \\
            }
        }
    }
    & & \\ \\

    \theutterance \stepcounter{utterance}  
    & & & \multicolumn{2}{p{0.3\linewidth}}{
        \cellcolor[rgb]{0.9,0.9,0.9}{
            \makecell[{{p{\linewidth}}}]{
                \texttt{\tiny{[GM$|$GM]}}
                \texttt{SÖYLE: Harika! O zaman ben ilk olarak en uzak nesnelerden biri olan D (5,7) nesnesini merkeze yaklaştırmak için hareket ettireceğim. Ben D nesnesini (4,4) konumuna taşımayı deniyorum. Sen de kendi tahtanda uygun bir nesneyi merkeze doğru hareket ettirebilirsin.} \\
            }
        }
    }
    & & \\ \\

    \theutterance \stepcounter{utterance}  
    & & \multicolumn{4}{p{0.6\linewidth}}{
        \cellcolor[rgb]{0.9,0.9,0.9}{
            \makecell[{{p{\linewidth}}}]{
                \texttt{\tiny{[GM$\rangle$P2]}}
                \texttt{Mesajın diğer oyuncuya iletildi.} \\
\\ 
\\ 
\texttt{Şu anda en fazla 28 turdan 3. turu oynuyorsunuz.} \\
\texttt{Toplamda 16 cezanın 1 kadarını biriktirdiniz.Diğer oyuncu şu mesajı gönderdi:} \\
\\ 
\\ 
\texttt{"Harika! O zaman ben ilk olarak en uzak nesnelerden biri olan D (5,7) nesnesini merkeze yaklaştırmak için hareket ettireceğim. Ben D nesnesini (4,4) konumuna taşımayı deniyorum. Sen de kendi tahtanda uygun bir nesneyi merkeze doğru hareket ettirebilirsin."} \\
\\ 
\\ 
\texttt{Bir sonraki komutun nedir?} \\
            }
        }
    }
    & \\ \\

    \theutterance \stepcounter{utterance}  
    & & & \multicolumn{4}{p{0.6\linewidth}}{
        \cellcolor[rgb]{0.9,0.9,0.9}{
            \makecell[{{p{\linewidth}}}]{
                \texttt{\tiny{[GM$\langle$P2]}}
                \texttt{HAMLE: D, (1,6)} \\
            }
        }
    }
     \\ \\

    \theutterance \stepcounter{utterance}  
    & & & \multicolumn{2}{p{0.3\linewidth}}{
        \cellcolor[rgb]{0.9,0.9,0.9}{
            \makecell[{{p{\linewidth}}}]{
                \texttt{\tiny{[GM$|$GM]}}
                \texttt{HAMLE: D, (1,6)} \\
            }
        }
    }
    & & \\ \\

    \theutterance \stepcounter{utterance}  
    & & & \multicolumn{2}{p{0.3\linewidth}}{
        \cellcolor[rgb]{0.9,0.9,0.9}{
            \makecell[{{p{\linewidth}}}]{
                \texttt{\tiny{[GM$|$GM]}}
                \texttt{Ceza: (1,6) boş değil; "─" içeriyor.} \\
\\ 
\texttt{Toplamda 16 cezanın 2 kadarını biriktirdiniz.Lütfen tekrar deneyin!} \\
            }
        }
    }
    & & \\ \\

    \theutterance \stepcounter{utterance}  
    & & \multicolumn{4}{p{0.6\linewidth}}{
        \cellcolor[rgb]{0.9,0.9,0.9}{
            \makecell[{{p{\linewidth}}}]{
                \texttt{\tiny{[GM$\rangle$P2]}}
                \texttt{Ceza: (1,6) boş değil; "─" içeriyor.} \\
\\ 
\texttt{Toplamda 16 cezanın 2 kadarını biriktirdiniz.Lütfen tekrar deneyin!} \\
            }
        }
    }
    & \\ \\

    \theutterance \stepcounter{utterance}  
    & & & \multicolumn{4}{p{0.6\linewidth}}{
        \cellcolor[rgb]{0.9,0.9,0.9}{
            \makecell[{{p{\linewidth}}}]{
                \texttt{\tiny{[GM$\langle$P2]}}
                \texttt{HAMLE: D, (6,7)} \\
            }
        }
    }
     \\ \\

    \theutterance \stepcounter{utterance}  
    & & & \multicolumn{2}{p{0.3\linewidth}}{
        \cellcolor[rgb]{0.9,0.9,0.9}{
            \makecell[{{p{\linewidth}}}]{
                \texttt{\tiny{[GM$|$GM]}}
                \texttt{HAMLE: D, (6,7)} \\
            }
        }
    }
    & & \\ \\

    \theutterance \stepcounter{utterance}  
    & & & \multicolumn{2}{p{0.3\linewidth}}{
        \cellcolor[rgb]{0.9,0.9,0.9}{
            \makecell[{{p{\linewidth}}}]{
                \texttt{\tiny{[GM$|$GM]}}
                \texttt{D moved to (6, 7)} \\
            }
        }
    }
    & & \\ \\

    \theutterance \stepcounter{utterance}  
    & & \multicolumn{4}{p{0.6\linewidth}}{
        \cellcolor[rgb]{0.9,0.9,0.9}{
            \makecell[{{p{\linewidth}}}]{
                \texttt{\tiny{[P1$\langle$GM]}}
                \texttt{Mesajın diğer oyuncuya iletildi.} \\
\\ 
\\ 
\texttt{Şu anda en fazla 28 turdan 3. turu oynuyorsunuz.} \\
\texttt{Toplamda 16 cezanın 2 kadarını biriktirdiniz.Diğer oyuncu kendi ızgarasında bir nesneyi hareket ettirdi.} \\
\\ 
\\ 
\texttt{Bir sonraki komutun nedir?} \\
            }
        }
    }
    & \\ \\

    \theutterance \stepcounter{utterance}  
    & \multicolumn{4}{p{0.6\linewidth}}{
        \cellcolor[rgb]{0.9,0.9,0.9}{
            \makecell[{{p{\linewidth}}}]{
                \texttt{\tiny{[P1$\rangle$GM]}}
                \texttt{HAMLE: D, (4, 4)} \\
            }
        }
    }
    & & \\ \\

    \theutterance \stepcounter{utterance}  
    & & & \multicolumn{2}{p{0.3\linewidth}}{
        \cellcolor[rgb]{0.9,0.9,0.9}{
            \makecell[{{p{\linewidth}}}]{
                \texttt{\tiny{[GM$|$GM]}}
                \texttt{HAMLE: D, (4, 4)} \\
            }
        }
    }
    & & \\ \\

    \theutterance \stepcounter{utterance}  
    & & & \multicolumn{2}{p{0.3\linewidth}}{
        \cellcolor[rgb]{0.9,0.9,0.9}{
            \makecell[{{p{\linewidth}}}]{
                \texttt{\tiny{[GM$|$GM]}}
                \texttt{Ceza: (4,4) boş değil; "─" içeriyor.} \\
\\ 
\texttt{Toplamda 16 cezanın 3 kadarını biriktirdiniz.Lütfen tekrar deneyin!} \\
            }
        }
    }
    & & \\ \\

    \theutterance \stepcounter{utterance}  
    & & \multicolumn{4}{p{0.6\linewidth}}{
        \cellcolor[rgb]{0.9,0.9,0.9}{
            \makecell[{{p{\linewidth}}}]{
                \texttt{\tiny{[P1$\langle$GM]}}
                \texttt{Ceza: (4,4) boş değil; "─" içeriyor.} \\
\\ 
\texttt{Toplamda 16 cezanın 3 kadarını biriktirdiniz.Lütfen tekrar deneyin!} \\
            }
        }
    }
    & \\ \\

    \theutterance \stepcounter{utterance}  
    & \multicolumn{4}{p{0.6\linewidth}}{
        \cellcolor[rgb]{0.9,0.9,0.9}{
            \makecell[{{p{\linewidth}}}]{
                \texttt{\tiny{[P1$\rangle$GM]}}
                \texttt{HAMLE: D, (5, 6)} \\
            }
        }
    }
    & & \\ \\

    \theutterance \stepcounter{utterance}  
    & & & \multicolumn{2}{p{0.3\linewidth}}{
        \cellcolor[rgb]{0.9,0.9,0.9}{
            \makecell[{{p{\linewidth}}}]{
                \texttt{\tiny{[GM$|$GM]}}
                \texttt{HAMLE: D, (5, 6)} \\
            }
        }
    }
    & & \\ \\

    \theutterance \stepcounter{utterance}  
    & & & \multicolumn{2}{p{0.3\linewidth}}{
        \cellcolor[rgb]{0.9,0.9,0.9}{
            \makecell[{{p{\linewidth}}}]{
                \texttt{\tiny{[GM$|$GM]}}
                \texttt{Ceza: (5,6) boş değil; "┤" içeriyor.} \\
\\ 
\texttt{Toplamda 16 cezanın 4 kadarını biriktirdiniz.Lütfen tekrar deneyin!} \\
            }
        }
    }
    & & \\ \\

    \theutterance \stepcounter{utterance}  
    & & \multicolumn{4}{p{0.6\linewidth}}{
        \cellcolor[rgb]{0.9,0.9,0.9}{
            \makecell[{{p{\linewidth}}}]{
                \texttt{\tiny{[P1$\langle$GM]}}
                \texttt{Ceza: (5,6) boş değil; "┤" içeriyor.} \\
\\ 
\texttt{Toplamda 16 cezanın 4 kadarını biriktirdiniz.Lütfen tekrar deneyin!} \\
            }
        }
    }
    & \\ \\

    \theutterance \stepcounter{utterance}  
    & \multicolumn{4}{p{0.6\linewidth}}{
        \cellcolor[rgb]{0.9,0.9,0.9}{
            \makecell[{{p{\linewidth}}}]{
                \texttt{\tiny{[P1$\rangle$GM]}}
                \texttt{HAMLE: D, (5, 5)} \\
            }
        }
    }
    & & \\ \\

    \theutterance \stepcounter{utterance}  
    & & & \multicolumn{2}{p{0.3\linewidth}}{
        \cellcolor[rgb]{0.9,0.9,0.9}{
            \makecell[{{p{\linewidth}}}]{
                \texttt{\tiny{[GM$|$GM]}}
                \texttt{HAMLE: D, (5, 5)} \\
            }
        }
    }
    & & \\ \\

    \theutterance \stepcounter{utterance}  
    & & & \multicolumn{2}{p{0.3\linewidth}}{
        \cellcolor[rgb]{0.9,0.9,0.9}{
            \makecell[{{p{\linewidth}}}]{
                \texttt{\tiny{[GM$|$GM]}}
                \texttt{Ceza: (5,5) boş değil; "│" içeriyor.} \\
\\ 
\texttt{Toplamda 16 cezanın 5 kadarını biriktirdiniz.Lütfen tekrar deneyin!} \\
            }
        }
    }
    & & \\ \\

    \theutterance \stepcounter{utterance}  
    & & \multicolumn{4}{p{0.6\linewidth}}{
        \cellcolor[rgb]{0.9,0.9,0.9}{
            \makecell[{{p{\linewidth}}}]{
                \texttt{\tiny{[P1$\langle$GM]}}
                \texttt{Ceza: (5,5) boş değil; "│" içeriyor.} \\
\\ 
\texttt{Toplamda 16 cezanın 5 kadarını biriktirdiniz.Lütfen tekrar deneyin!} \\
            }
        }
    }
    & \\ \\

    \theutterance \stepcounter{utterance}  
    & \multicolumn{4}{p{0.6\linewidth}}{
        \cellcolor[rgb]{0.9,0.9,0.9}{
            \makecell[{{p{\linewidth}}}]{
                \texttt{\tiny{[P1$\rangle$GM]}}
                \texttt{HAMLE: D, (6, 7)} \\
            }
        }
    }
    & & \\ \\

    \theutterance \stepcounter{utterance}  
    & & & \multicolumn{2}{p{0.3\linewidth}}{
        \cellcolor[rgb]{0.9,0.9,0.9}{
            \makecell[{{p{\linewidth}}}]{
                \texttt{\tiny{[GM$|$GM]}}
                \texttt{HAMLE: D, (6, 7)} \\
            }
        }
    }
    & & \\ \\

    \theutterance \stepcounter{utterance}  
    & & & \multicolumn{2}{p{0.3\linewidth}}{
        \cellcolor[rgb]{0.9,0.9,0.9}{
            \makecell[{{p{\linewidth}}}]{
                \texttt{\tiny{[GM$|$GM]}}
                \texttt{Ceza: (6,7) boş değil; "D" içeriyor.} \\
\\ 
\texttt{Toplamda 16 cezanın 6 kadarını biriktirdiniz.Lütfen tekrar deneyin!} \\
            }
        }
    }
    & & \\ \\

    \theutterance \stepcounter{utterance}  
    & & \multicolumn{4}{p{0.6\linewidth}}{
        \cellcolor[rgb]{0.9,0.9,0.9}{
            \makecell[{{p{\linewidth}}}]{
                \texttt{\tiny{[P1$\langle$GM]}}
                \texttt{Ceza: (6,7) boş değil; "D" içeriyor.} \\
\\ 
\texttt{Toplamda 16 cezanın 6 kadarını biriktirdiniz.Lütfen tekrar deneyin!} \\
            }
        }
    }
    & \\ \\

    \theutterance \stepcounter{utterance}  
    & \multicolumn{4}{p{0.6\linewidth}}{
        \cellcolor[rgb]{0.9,0.9,0.9}{
            \makecell[{{p{\linewidth}}}]{
                \texttt{\tiny{[P1$\rangle$GM]}}
                \texttt{HAMLE: D, (6, 5)} \\
            }
        }
    }
    & & \\ \\

    \theutterance \stepcounter{utterance}  
    & & & \multicolumn{2}{p{0.3\linewidth}}{
        \cellcolor[rgb]{0.9,0.9,0.9}{
            \makecell[{{p{\linewidth}}}]{
                \texttt{\tiny{[GM$|$GM]}}
                \texttt{HAMLE: D, (6, 5)} \\
            }
        }
    }
    & & \\ \\

    \theutterance \stepcounter{utterance}  
    & & & \multicolumn{2}{p{0.3\linewidth}}{
        \cellcolor[rgb]{0.9,0.9,0.9}{
            \makecell[{{p{\linewidth}}}]{
                \texttt{\tiny{[GM$|$GM]}}
                \texttt{D moved to (6, 5)} \\
            }
        }
    }
    & & \\ \\

    \theutterance \stepcounter{utterance}  
    & & \multicolumn{4}{p{0.6\linewidth}}{
        \cellcolor[rgb]{0.9,0.9,0.9}{
            \makecell[{{p{\linewidth}}}]{
                \texttt{\tiny{[GM$\rangle$P2]}}
                \texttt{D’ başarıyla (6,7) koordinatlarına taşındı. Güncellenmiş ızgaran şu şekilde:} \\
\texttt{```} \\
\texttt{ 1234567} \\
\texttt{╔════╤══╗} \\
\texttt{║◌◌◌◌│O◌║ 1} \\
\texttt{║P◌M◌├──╢ 2} \\
\texttt{║◌S◌◌│T◌║ 3} \\
\texttt{║◌┌──┼──╢ 4} \\
\texttt{║A│◌◌│◌◌║ 5} \\
\texttt{╟─┴┬─┤◌◌║ 6} \\
\texttt{║◌◌│◌│D◌║ 7} \\
\texttt{╚══╧═╧══╝} \\
\texttt{```} \\
\texttt{Şu anda en fazla 28 turdan 4. turu oynuyorsunuz.} \\
\texttt{Toplamda 16 cezanın 6 kadarını biriktirdiniz.Diğer oyuncu kendi ızgarasında bir nesneyi hareket ettirdi.} \\
\\ 
\\ 
\texttt{Bir sonraki komutun nedir?} \\
            }
        }
    }
    & \\ \\

    \theutterance \stepcounter{utterance}  
    & & & \multicolumn{4}{p{0.6\linewidth}}{
        \cellcolor[rgb]{0.9,0.9,0.9}{
            \makecell[{{p{\linewidth}}}]{
                \texttt{\tiny{[GM$\langle$P2]}}
                \texttt{HAMLE: A, (1,4)} \\
            }
        }
    }
     \\ \\

    \theutterance \stepcounter{utterance}  
    & & & \multicolumn{2}{p{0.3\linewidth}}{
        \cellcolor[rgb]{0.9,0.9,0.9}{
            \makecell[{{p{\linewidth}}}]{
                \texttt{\tiny{[GM$|$GM]}}
                \texttt{HAMLE: A, (1,4)} \\
            }
        }
    }
    & & \\ \\

    \theutterance \stepcounter{utterance}  
    & & & \multicolumn{2}{p{0.3\linewidth}}{
        \cellcolor[rgb]{0.9,0.9,0.9}{
            \makecell[{{p{\linewidth}}}]{
                \texttt{\tiny{[GM$|$GM]}}
                \texttt{A moved to (1, 4)} \\
            }
        }
    }
    & & \\ \\

    \theutterance \stepcounter{utterance}  
    & & \multicolumn{4}{p{0.6\linewidth}}{
        \cellcolor[rgb]{0.9,0.9,0.9}{
            \makecell[{{p{\linewidth}}}]{
                \texttt{\tiny{[P1$\langle$GM]}}
                \texttt{D’ başarıyla (6,5) koordinatlarına taşındı. Güncellenmiş ızgaran şu şekilde:} \\
\texttt{```} \\
\texttt{ 1234567} \\
\texttt{╔════╤══╗} \\
\texttt{║◌◌◌O│PM║ 1} \\
\texttt{║◌◌T◌├──╢ 2} \\
\texttt{║◌◌◌◌│◌◌║ 3} \\
\texttt{║S┌──┼──╢ 4} \\
\texttt{║◌│◌A│D◌║ 5} \\
\texttt{╟─┴┬─┤◌◌║ 6} \\
\texttt{║◌◌│◌│◌◌║ 7} \\
\texttt{╚══╧═╧══╝} \\
\texttt{```} \\
\texttt{Şu anda en fazla 28 turdan 4. turu oynuyorsunuz.} \\
\texttt{Toplamda 16 cezanın 6 kadarını biriktirdiniz.Diğer oyuncu kendi ızgarasında bir nesneyi hareket ettirdi.} \\
\\ 
\\ 
\texttt{Bir sonraki komutun nedir?} \\
            }
        }
    }
    & \\ \\

    \theutterance \stepcounter{utterance}  
    & \multicolumn{4}{p{0.6\linewidth}}{
        \cellcolor[rgb]{0.9,0.9,0.9}{
            \makecell[{{p{\linewidth}}}]{
                \texttt{\tiny{[P1$\rangle$GM]}}
                \texttt{SÖYLE: D nesnesini (6,5) konumuna başarıyla taşıdım. Diğer oyuncu hangi nesneyi nereye hareket ettirdi? Bir sonraki hamlem için plan yapmamıza yardımcı olması adına onun hamlesini de paylaşabilir misin?} \\
            }
        }
    }
    & & \\ \\

    \theutterance \stepcounter{utterance}  
    & & & \multicolumn{2}{p{0.3\linewidth}}{
        \cellcolor[rgb]{0.9,0.9,0.9}{
            \makecell[{{p{\linewidth}}}]{
                \texttt{\tiny{[GM$|$GM]}}
                \texttt{SÖYLE: D nesnesini (6,5) konumuna başarıyla taşıdım. Diğer oyuncu hangi nesneyi nereye hareket ettirdi? Bir sonraki hamlem için plan yapmamıza yardımcı olması adına onun hamlesini de paylaşabilir misin?} \\
            }
        }
    }
    & & \\ \\

    \theutterance \stepcounter{utterance}  
    & & \multicolumn{4}{p{0.6\linewidth}}{
        \cellcolor[rgb]{0.9,0.9,0.9}{
            \makecell[{{p{\linewidth}}}]{
                \texttt{\tiny{[GM$\rangle$P2]}}
                \texttt{A’ başarıyla (1,4) koordinatlarına taşındı. Güncellenmiş ızgaran şu şekilde:} \\
\texttt{```} \\
\texttt{ 1234567} \\
\texttt{╔════╤══╗} \\
\texttt{║◌◌◌◌│O◌║ 1} \\
\texttt{║P◌M◌├──╢ 2} \\
\texttt{║◌S◌◌│T◌║ 3} \\
\texttt{║A┌──┼──╢ 4} \\
\texttt{║◌│◌◌│◌◌║ 5} \\
\texttt{╟─┴┬─┤◌◌║ 6} \\
\texttt{║◌◌│◌│D◌║ 7} \\
\texttt{╚══╧═╧══╝} \\
\texttt{```} \\
\texttt{Şu anda en fazla 28 turdan 5. turu oynuyorsunuz.} \\
\texttt{Toplamda 16 cezanın 6 kadarını biriktirdiniz.Diğer oyuncu şu mesajı gönderdi:} \\
\\ 
\\ 
\texttt{"D nesnesini (6,5) konumuna başarıyla taşıdım. Diğer oyuncu hangi nesneyi nereye hareket ettirdi? Bir sonraki hamlem için plan yapmamıza yardımcı olması adına onun hamlesini de paylaşabilir misin?"} \\
\\ 
\\ 
\texttt{Bir sonraki komutun nedir?} \\
            }
        }
    }
    & \\ \\

    \theutterance \stepcounter{utterance}  
    & & & \multicolumn{4}{p{0.6\linewidth}}{
        \cellcolor[rgb]{0.9,0.9,0.9}{
            \makecell[{{p{\linewidth}}}]{
                \texttt{\tiny{[GM$\langle$P2]}}
                \texttt{SÖYLE: Ben A nesnesini (1,4) konumuna taşıdım. Şimdi (1,1) civarında bir küme oluşturmaya devam edebiliriz. Bir sonraki hamlemde S nesnesini merkeze yaklaştırmayı düşünüyorum.} \\
            }
        }
    }
     \\ \\

    \theutterance \stepcounter{utterance}  
    & & & \multicolumn{2}{p{0.3\linewidth}}{
        \cellcolor[rgb]{0.9,0.9,0.9}{
            \makecell[{{p{\linewidth}}}]{
                \texttt{\tiny{[GM$|$GM]}}
                \texttt{SÖYLE: Ben A nesnesini (1,4) konumuna taşıdım. Şimdi (1,1) civarında bir küme oluşturmaya devam edebiliriz. Bir sonraki hamlemde S nesnesini merkeze yaklaştırmayı düşünüyorum.} \\
            }
        }
    }
    & & \\ \\

    \theutterance \stepcounter{utterance}  
    & & \multicolumn{4}{p{0.6\linewidth}}{
        \cellcolor[rgb]{0.9,0.9,0.9}{
            \makecell[{{p{\linewidth}}}]{
                \texttt{\tiny{[P1$\langle$GM]}}
                \texttt{Mesajın diğer oyuncuya iletildi.} \\
\\ 
\\ 
\texttt{Şu anda en fazla 28 turdan 5. turu oynuyorsunuz.} \\
\texttt{Toplamda 16 cezanın 6 kadarını biriktirdiniz.Diğer oyuncu şu mesajı gönderdi:} \\
\\ 
\\ 
\texttt{"Ben A nesnesini (1,4) konumuna taşıdım. Şimdi (1,1) civarında bir küme oluşturmaya devam edebiliriz. Bir sonraki hamlemde S nesnesini merkeze yaklaştırmayı düşünüyorum."} \\
\\ 
\\ 
\texttt{Bir sonraki komutun nedir?} \\
            }
        }
    }
    & \\ \\

    \theutterance \stepcounter{utterance}  
    & \multicolumn{4}{p{0.6\linewidth}}{
        \cellcolor[rgb]{0.9,0.9,0.9}{
            \makecell[{{p{\linewidth}}}]{
                \texttt{\tiny{[P1$\rangle$GM]}}
                \texttt{SÖYLE: Anladım, A nesnesini (1,4) konumuna taşıman güzel bir adım. Ben de kendi tahtamdaki nesneleri merkeze yaklaştırmaya devam edeceğim. Bir sonraki hamlemde S nesnesini (1,4) konumuna taşımayı deneyeceğim.} \\
            }
        }
    }
    & & \\ \\

    \theutterance \stepcounter{utterance}  
    & & & \multicolumn{2}{p{0.3\linewidth}}{
        \cellcolor[rgb]{0.9,0.9,0.9}{
            \makecell[{{p{\linewidth}}}]{
                \texttt{\tiny{[GM$|$GM]}}
                \texttt{SÖYLE: Anladım, A nesnesini (1,4) konumuna taşıman güzel bir adım. Ben de kendi tahtamdaki nesneleri merkeze yaklaştırmaya devam edeceğim. Bir sonraki hamlemde S nesnesini (1,4) konumuna taşımayı deneyeceğim.} \\
            }
        }
    }
    & & \\ \\

    \theutterance \stepcounter{utterance}  
    & & \multicolumn{4}{p{0.6\linewidth}}{
        \cellcolor[rgb]{0.9,0.9,0.9}{
            \makecell[{{p{\linewidth}}}]{
                \texttt{\tiny{[GM$\rangle$P2]}}
                \texttt{Mesajın diğer oyuncuya iletildi.} \\
\\ 
\\ 
\texttt{Şu anda en fazla 28 turdan 6. turu oynuyorsunuz.} \\
\texttt{Toplamda 16 cezanın 6 kadarını biriktirdiniz.Diğer oyuncu şu mesajı gönderdi:} \\
\\ 
\\ 
\texttt{"Anladım, A nesnesini (1,4) konumuna taşıman güzel bir adım. Ben de kendi tahtamdaki nesneleri merkeze yaklaştırmaya devam edeceğim. Bir sonraki hamlemde S nesnesini (1,4) konumuna taşımayı deneyeceğim."} \\
\\ 
\\ 
\texttt{Bir sonraki komutun nedir?} \\
            }
        }
    }
    & \\ \\

    \theutterance \stepcounter{utterance}  
    & & & \multicolumn{4}{p{0.6\linewidth}}{
        \cellcolor[rgb]{0.9,0.9,0.9}{
            \makecell[{{p{\linewidth}}}]{
                \texttt{\tiny{[GM$\langle$P2]}}
                \texttt{HAMLE: S, (1,2)} \\
            }
        }
    }
     \\ \\

    \theutterance \stepcounter{utterance}  
    & & & \multicolumn{2}{p{0.3\linewidth}}{
        \cellcolor[rgb]{0.9,0.9,0.9}{
            \makecell[{{p{\linewidth}}}]{
                \texttt{\tiny{[GM$|$GM]}}
                \texttt{HAMLE: S, (1,2)} \\
            }
        }
    }
    & & \\ \\

    \theutterance \stepcounter{utterance}  
    & & & \multicolumn{2}{p{0.3\linewidth}}{
        \cellcolor[rgb]{0.9,0.9,0.9}{
            \makecell[{{p{\linewidth}}}]{
                \texttt{\tiny{[GM$|$GM]}}
                \texttt{Ceza: (1,2) boş değil; "P" içeriyor.} \\
\\ 
\texttt{Toplamda 16 cezanın 7 kadarını biriktirdiniz.Lütfen tekrar deneyin!} \\
            }
        }
    }
    & & \\ \\

    \theutterance \stepcounter{utterance}  
    & & \multicolumn{4}{p{0.6\linewidth}}{
        \cellcolor[rgb]{0.9,0.9,0.9}{
            \makecell[{{p{\linewidth}}}]{
                \texttt{\tiny{[GM$\rangle$P2]}}
                \texttt{Ceza: (1,2) boş değil; "P" içeriyor.} \\
\\ 
\texttt{Toplamda 16 cezanın 7 kadarını biriktirdiniz.Lütfen tekrar deneyin!} \\
            }
        }
    }
    & \\ \\

    \theutterance \stepcounter{utterance}  
    & & & \multicolumn{4}{p{0.6\linewidth}}{
        \cellcolor[rgb]{0.9,0.9,0.9}{
            \makecell[{{p{\linewidth}}}]{
                \texttt{\tiny{[GM$\langle$P2]}}
                \texttt{HAMLE: S, (2,3)} \\
            }
        }
    }
     \\ \\

    \theutterance \stepcounter{utterance}  
    & & & \multicolumn{2}{p{0.3\linewidth}}{
        \cellcolor[rgb]{0.9,0.9,0.9}{
            \makecell[{{p{\linewidth}}}]{
                \texttt{\tiny{[GM$|$GM]}}
                \texttt{HAMLE: S, (2,3)} \\
            }
        }
    }
    & & \\ \\

    \theutterance \stepcounter{utterance}  
    & & & \multicolumn{2}{p{0.3\linewidth}}{
        \cellcolor[rgb]{0.9,0.9,0.9}{
            \makecell[{{p{\linewidth}}}]{
                \texttt{\tiny{[GM$|$GM]}}
                \texttt{Ceza: (2,3) boş değil; "S" içeriyor.} \\
\\ 
\texttt{Toplamda 16 cezanın 8 kadarını biriktirdiniz.Lütfen tekrar deneyin!} \\
            }
        }
    }
    & & \\ \\

    \theutterance \stepcounter{utterance}  
    & & \multicolumn{4}{p{0.6\linewidth}}{
        \cellcolor[rgb]{0.9,0.9,0.9}{
            \makecell[{{p{\linewidth}}}]{
                \texttt{\tiny{[GM$\rangle$P2]}}
                \texttt{Ceza: (2,3) boş değil; "S" içeriyor.} \\
\\ 
\texttt{Toplamda 16 cezanın 8 kadarını biriktirdiniz.Lütfen tekrar deneyin!} \\
            }
        }
    }
    & \\ \\

    \theutterance \stepcounter{utterance}  
    & & & \multicolumn{4}{p{0.6\linewidth}}{
        \cellcolor[rgb]{0.9,0.9,0.9}{
            \makecell[{{p{\linewidth}}}]{
                \texttt{\tiny{[GM$\langle$P2]}}
                \texttt{HAMLE: S, (2,4)} \\
            }
        }
    }
     \\ \\

    \theutterance \stepcounter{utterance}  
    & & & \multicolumn{2}{p{0.3\linewidth}}{
        \cellcolor[rgb]{0.9,0.9,0.9}{
            \makecell[{{p{\linewidth}}}]{
                \texttt{\tiny{[GM$|$GM]}}
                \texttt{HAMLE: S, (2,4)} \\
            }
        }
    }
    & & \\ \\

    \theutterance \stepcounter{utterance}  
    & & & \multicolumn{2}{p{0.3\linewidth}}{
        \cellcolor[rgb]{0.9,0.9,0.9}{
            \makecell[{{p{\linewidth}}}]{
                \texttt{\tiny{[GM$|$GM]}}
                \texttt{Ceza: (2,4) boş değil; "┌" içeriyor.} \\
\\ 
\texttt{Toplamda 16 cezanın 9 kadarını biriktirdiniz.Lütfen tekrar deneyin!} \\
            }
        }
    }
    & & \\ \\

    \theutterance \stepcounter{utterance}  
    & & \multicolumn{4}{p{0.6\linewidth}}{
        \cellcolor[rgb]{0.9,0.9,0.9}{
            \makecell[{{p{\linewidth}}}]{
                \texttt{\tiny{[GM$\rangle$P2]}}
                \texttt{Ceza: (2,4) boş değil; "┌" içeriyor.} \\
\\ 
\texttt{Toplamda 16 cezanın 9 kadarını biriktirdiniz.Lütfen tekrar deneyin!} \\
            }
        }
    }
    & \\ \\

    \theutterance \stepcounter{utterance}  
    & & & \multicolumn{4}{p{0.6\linewidth}}{
        \cellcolor[rgb]{0.9,0.9,0.9}{
            \makecell[{{p{\linewidth}}}]{
                \texttt{\tiny{[GM$\langle$P2]}}
                \texttt{HAMLE: S, (2,2)} \\
            }
        }
    }
     \\ \\

    \theutterance \stepcounter{utterance}  
    & & & \multicolumn{2}{p{0.3\linewidth}}{
        \cellcolor[rgb]{0.9,0.9,0.9}{
            \makecell[{{p{\linewidth}}}]{
                \texttt{\tiny{[GM$|$GM]}}
                \texttt{HAMLE: S, (2,2)} \\
            }
        }
    }
    & & \\ \\

    \theutterance \stepcounter{utterance}  
    & & & \multicolumn{2}{p{0.3\linewidth}}{
        \cellcolor[rgb]{0.9,0.9,0.9}{
            \makecell[{{p{\linewidth}}}]{
                \texttt{\tiny{[GM$|$GM]}}
                \texttt{S moved to (2, 2)} \\
            }
        }
    }
    & & \\ \\

    \theutterance \stepcounter{utterance}  
    & & \multicolumn{4}{p{0.6\linewidth}}{
        \cellcolor[rgb]{0.9,0.9,0.9}{
            \makecell[{{p{\linewidth}}}]{
                \texttt{\tiny{[P1$\langle$GM]}}
                \texttt{Mesajın diğer oyuncuya iletildi.} \\
\\ 
\\ 
\texttt{Şu anda en fazla 28 turdan 6. turu oynuyorsunuz.} \\
\texttt{Toplamda 16 cezanın 9 kadarını biriktirdiniz.Diğer oyuncu kendi ızgarasında bir nesneyi hareket ettirdi.} \\
\\ 
\\ 
\texttt{Bir sonraki komutun nedir?} \\
            }
        }
    }
    & \\ \\

    \theutterance \stepcounter{utterance}  
    & \multicolumn{4}{p{0.6\linewidth}}{
        \cellcolor[rgb]{0.9,0.9,0.9}{
            \makecell[{{p{\linewidth}}}]{
                \texttt{\tiny{[P1$\rangle$GM]}}
                \texttt{HAMLE: S, (1, 4)} \\
            }
        }
    }
    & & \\ \\

    \theutterance \stepcounter{utterance}  
    & & & \multicolumn{2}{p{0.3\linewidth}}{
        \cellcolor[rgb]{0.9,0.9,0.9}{
            \makecell[{{p{\linewidth}}}]{
                \texttt{\tiny{[GM$|$GM]}}
                \texttt{HAMLE: S, (1, 4)} \\
            }
        }
    }
    & & \\ \\

    \theutterance \stepcounter{utterance}  
    & & & \multicolumn{2}{p{0.3\linewidth}}{
        \cellcolor[rgb]{0.9,0.9,0.9}{
            \makecell[{{p{\linewidth}}}]{
                \texttt{\tiny{[GM$|$GM]}}
                \texttt{Ceza: (1,4) boş değil; "S" içeriyor.} \\
\\ 
\texttt{Toplamda 16 cezanın 10 kadarını biriktirdiniz.Lütfen tekrar deneyin!} \\
            }
        }
    }
    & & \\ \\

    \theutterance \stepcounter{utterance}  
    & & \multicolumn{4}{p{0.6\linewidth}}{
        \cellcolor[rgb]{0.9,0.9,0.9}{
            \makecell[{{p{\linewidth}}}]{
                \texttt{\tiny{[P1$\langle$GM]}}
                \texttt{Ceza: (1,4) boş değil; "S" içeriyor.} \\
\\ 
\texttt{Toplamda 16 cezanın 10 kadarını biriktirdiniz.Lütfen tekrar deneyin!} \\
            }
        }
    }
    & \\ \\

    \theutterance \stepcounter{utterance}  
    & \multicolumn{4}{p{0.6\linewidth}}{
        \cellcolor[rgb]{0.9,0.9,0.9}{
            \makecell[{{p{\linewidth}}}]{
                \texttt{\tiny{[P1$\rangle$GM]}}
                \texttt{HAMLE: S, (1, 3)} \\
            }
        }
    }
    & & \\ \\

    \theutterance \stepcounter{utterance}  
    & & & \multicolumn{2}{p{0.3\linewidth}}{
        \cellcolor[rgb]{0.9,0.9,0.9}{
            \makecell[{{p{\linewidth}}}]{
                \texttt{\tiny{[GM$|$GM]}}
                \texttt{HAMLE: S, (1, 3)} \\
            }
        }
    }
    & & \\ \\

    \theutterance \stepcounter{utterance}  
    & & & \multicolumn{2}{p{0.3\linewidth}}{
        \cellcolor[rgb]{0.9,0.9,0.9}{
            \makecell[{{p{\linewidth}}}]{
                \texttt{\tiny{[GM$|$GM]}}
                \texttt{S moved to (1, 3)} \\
            }
        }
    }
    & & \\ \\

    \theutterance \stepcounter{utterance}  
    & & \multicolumn{4}{p{0.6\linewidth}}{
        \cellcolor[rgb]{0.9,0.9,0.9}{
            \makecell[{{p{\linewidth}}}]{
                \texttt{\tiny{[GM$\rangle$P2]}}
                \texttt{S’ başarıyla (2,2) koordinatlarına taşındı. Güncellenmiş ızgaran şu şekilde:} \\
\texttt{```} \\
\texttt{ 1234567} \\
\texttt{╔════╤══╗} \\
\texttt{║◌◌◌◌│O◌║ 1} \\
\texttt{║PSM◌├──╢ 2} \\
\texttt{║◌◌◌◌│T◌║ 3} \\
\texttt{║A┌──┼──╢ 4} \\
\texttt{║◌│◌◌│◌◌║ 5} \\
\texttt{╟─┴┬─┤◌◌║ 6} \\
\texttt{║◌◌│◌│D◌║ 7} \\
\texttt{╚══╧═╧══╝} \\
\texttt{```} \\
\texttt{Şu anda en fazla 28 turdan 7. turu oynuyorsunuz.} \\
\texttt{Toplamda 16 cezanın 10 kadarını biriktirdiniz.Diğer oyuncu kendi ızgarasında bir nesneyi hareket ettirdi.} \\
\\ 
\\ 
\texttt{Bir sonraki komutun nedir?} \\
            }
        }
    }
    & \\ \\

    \theutterance \stepcounter{utterance}  
    & & & \multicolumn{4}{p{0.6\linewidth}}{
        \cellcolor[rgb]{0.9,0.9,0.9}{
            \makecell[{{p{\linewidth}}}]{
                \texttt{\tiny{[GM$\langle$P2]}}
                \texttt{SÖYLE: S nesnesini (2,2) konumuna başarıyla taşıdım. Şimdi (1,1) civarındaki kümemizi tamamlamak için P veya M nesnelerini hareket ettirebiliriz. Sen kendi tahtanda hangi nesneyi nereye taşıdın?} \\
            }
        }
    }
     \\ \\

    \theutterance \stepcounter{utterance}  
    & & & \multicolumn{2}{p{0.3\linewidth}}{
        \cellcolor[rgb]{0.9,0.9,0.9}{
            \makecell[{{p{\linewidth}}}]{
                \texttt{\tiny{[GM$|$GM]}}
                \texttt{SÖYLE: S nesnesini (2,2) konumuna başarıyla taşıdım. Şimdi (1,1) civarındaki kümemizi tamamlamak için P veya M nesnelerini hareket ettirebiliriz. Sen kendi tahtanda hangi nesneyi nereye taşıdın?} \\
            }
        }
    }
    & & \\ \\

    \theutterance \stepcounter{utterance}  
    & & \multicolumn{4}{p{0.6\linewidth}}{
        \cellcolor[rgb]{0.9,0.9,0.9}{
            \makecell[{{p{\linewidth}}}]{
                \texttt{\tiny{[P1$\langle$GM]}}
                \texttt{S’ başarıyla (1,3) koordinatlarına taşındı. Güncellenmiş ızgaran şu şekilde:} \\
\texttt{```} \\
\texttt{ 1234567} \\
\texttt{╔════╤══╗} \\
\texttt{║◌◌◌O│PM║ 1} \\
\texttt{║◌◌T◌├──╢ 2} \\
\texttt{║S◌◌◌│◌◌║ 3} \\
\texttt{║◌┌──┼──╢ 4} \\
\texttt{║◌│◌A│D◌║ 5} \\
\texttt{╟─┴┬─┤◌◌║ 6} \\
\texttt{║◌◌│◌│◌◌║ 7} \\
\texttt{╚══╧═╧══╝} \\
\texttt{```} \\
\texttt{Şu anda en fazla 28 turdan 7. turu oynuyorsunuz.} \\
\texttt{Toplamda 16 cezanın 10 kadarını biriktirdiniz.Diğer oyuncu şu mesajı gönderdi:} \\
\\ 
\\ 
\texttt{"S nesnesini (2,2) konumuna başarıyla taşıdım. Şimdi (1,1) civarındaki kümemizi tamamlamak için P veya M nesnelerini hareket ettirebiliriz. Sen kendi tahtanda hangi nesneyi nereye taşıdın?"} \\
\\ 
\\ 
\texttt{Bir sonraki komutun nedir?} \\
            }
        }
    }
    & \\ \\

    \theutterance \stepcounter{utterance}  
    & \multicolumn{4}{p{0.6\linewidth}}{
        \cellcolor[rgb]{0.9,0.9,0.9}{
            \makecell[{{p{\linewidth}}}]{
                \texttt{\tiny{[P1$\rangle$GM]}}
                \texttt{SÖYLE: Ben S nesnesini (1,3) konumuna taşıdım. P ve M nesneleri zaten (5,1) ve (6,1) konumlarında, yani üst kenarda duruyorlar. Onları merkeze yaklaştırmak için P nesnesini (4,1) konumuna, M nesnesini ise (5,1) konumuna çekebiliriz. Bir sonraki hamlemde P nesnesini hareket ettireceğim.} \\
            }
        }
    }
    & & \\ \\

    \theutterance \stepcounter{utterance}  
    & & & \multicolumn{2}{p{0.3\linewidth}}{
        \cellcolor[rgb]{0.9,0.9,0.9}{
            \makecell[{{p{\linewidth}}}]{
                \texttt{\tiny{[GM$|$GM]}}
                \texttt{SÖYLE: Ben S nesnesini (1,3) konumuna taşıdım. P ve M nesneleri zaten (5,1) ve (6,1) konumlarında, yani üst kenarda duruyorlar. Onları merkeze yaklaştırmak için P nesnesini (4,1) konumuna, M nesnesini ise (5,1) konumuna çekebiliriz. Bir sonraki hamlemde P nesnesini hareket ettireceğim.} \\
            }
        }
    }
    & & \\ \\

    \theutterance \stepcounter{utterance}  
    & & \multicolumn{4}{p{0.6\linewidth}}{
        \cellcolor[rgb]{0.9,0.9,0.9}{
            \makecell[{{p{\linewidth}}}]{
                \texttt{\tiny{[GM$\rangle$P2]}}
                \texttt{Mesajın diğer oyuncuya iletildi.} \\
\\ 
\\ 
\texttt{Şu anda en fazla 28 turdan 8. turu oynuyorsunuz.} \\
\texttt{Toplamda 16 cezanın 10 kadarını biriktirdiniz.Diğer oyuncu şu mesajı gönderdi:} \\
\\ 
\\ 
\texttt{"Ben S nesnesini (1,3) konumuna taşıdım. P ve M nesneleri zaten (5,1) ve (6,1) konumlarında, yani üst kenarda duruyorlar. Onları merkeze yaklaştırmak için P nesnesini (4,1) konumuna, M nesnesini ise (5,1) konumuna çekebiliriz. Bir sonraki hamlemde P nesnesini hareket ettireceğim."} \\
\\ 
\\ 
\texttt{Bir sonraki komutun nedir?} \\
            }
        }
    }
    & \\ \\

    \theutterance \stepcounter{utterance}  
    & & & \multicolumn{4}{p{0.6\linewidth}}{
        \cellcolor[rgb]{0.9,0.9,0.9}{
            \makecell[{{p{\linewidth}}}]{
                \texttt{\tiny{[GM$\langle$P2]}}
                \texttt{SÖYLE: P ve M nesnelerini merkeze çekme planın mantıklı. Ben de kendi tahtamdaki P (1,2) ve M (3,2) nesnelerini merkeze yaklaştırmaya başlayacağım. Bir sonraki hamlemde P nesnesini hareket ettireceğim.} \\
            }
        }
    }
     \\ \\

    \theutterance \stepcounter{utterance}  
    & & & \multicolumn{2}{p{0.3\linewidth}}{
        \cellcolor[rgb]{0.9,0.9,0.9}{
            \makecell[{{p{\linewidth}}}]{
                \texttt{\tiny{[GM$|$GM]}}
                \texttt{SÖYLE: P ve M nesnelerini merkeze çekme planın mantıklı. Ben de kendi tahtamdaki P (1,2) ve M (3,2) nesnelerini merkeze yaklaştırmaya başlayacağım. Bir sonraki hamlemde P nesnesini hareket ettireceğim.} \\
            }
        }
    }
    & & \\ \\

    \theutterance \stepcounter{utterance}  
    & & \multicolumn{4}{p{0.6\linewidth}}{
        \cellcolor[rgb]{0.9,0.9,0.9}{
            \makecell[{{p{\linewidth}}}]{
                \texttt{\tiny{[P1$\langle$GM]}}
                \texttt{Mesajın diğer oyuncuya iletildi.} \\
\\ 
\\ 
\texttt{Şu anda en fazla 28 turdan 8. turu oynuyorsunuz.} \\
\texttt{Toplamda 16 cezanın 10 kadarını biriktirdiniz.Diğer oyuncu şu mesajı gönderdi:} \\
\\ 
\\ 
\texttt{"P ve M nesnelerini merkeze çekme planın mantıklı. Ben de kendi tahtamdaki P (1,2) ve M (3,2) nesnelerini merkeze yaklaştırmaya başlayacağım. Bir sonraki hamlemde P nesnesini hareket ettireceğim."} \\
\\ 
\\ 
\texttt{Bir sonraki komutun nedir?} \\
            }
        }
    }
    & \\ \\

    \theutterance \stepcounter{utterance}  
    & \multicolumn{4}{p{0.6\linewidth}}{
        \cellcolor[rgb]{0.9,0.9,0.9}{
            \makecell[{{p{\linewidth}}}]{
                \texttt{\tiny{[P1$\rangle$GM]}}
                \texttt{HAMLE: P, (4, 1)} \\
            }
        }
    }
    & & \\ \\

    \theutterance \stepcounter{utterance}  
    & & & \multicolumn{2}{p{0.3\linewidth}}{
        \cellcolor[rgb]{0.9,0.9,0.9}{
            \makecell[{{p{\linewidth}}}]{
                \texttt{\tiny{[GM$|$GM]}}
                \texttt{HAMLE: P, (4, 1)} \\
            }
        }
    }
    & & \\ \\

    \theutterance \stepcounter{utterance}  
    & & & \multicolumn{2}{p{0.3\linewidth}}{
        \cellcolor[rgb]{0.9,0.9,0.9}{
            \makecell[{{p{\linewidth}}}]{
                \texttt{\tiny{[GM$|$GM]}}
                \texttt{Ceza: (4,1) boş değil; "O" içeriyor.} \\
\\ 
\texttt{Toplamda 16 cezanın 11 kadarını biriktirdiniz.Lütfen tekrar deneyin!} \\
            }
        }
    }
    & & \\ \\

    \theutterance \stepcounter{utterance}  
    & & \multicolumn{4}{p{0.6\linewidth}}{
        \cellcolor[rgb]{0.9,0.9,0.9}{
            \makecell[{{p{\linewidth}}}]{
                \texttt{\tiny{[P1$\langle$GM]}}
                \texttt{Ceza: (4,1) boş değil; "O" içeriyor.} \\
\\ 
\texttt{Toplamda 16 cezanın 11 kadarını biriktirdiniz.Lütfen tekrar deneyin!} \\
            }
        }
    }
    & \\ \\

    \theutterance \stepcounter{utterance}  
    & \multicolumn{4}{p{0.6\linewidth}}{
        \cellcolor[rgb]{0.9,0.9,0.9}{
            \makecell[{{p{\linewidth}}}]{
                \texttt{\tiny{[P1$\rangle$GM]}}
                \texttt{HAMLE: P, (5, 1)} \\
            }
        }
    }
    & & \\ \\

    \theutterance \stepcounter{utterance}  
    & & & \multicolumn{2}{p{0.3\linewidth}}{
        \cellcolor[rgb]{0.9,0.9,0.9}{
            \makecell[{{p{\linewidth}}}]{
                \texttt{\tiny{[GM$|$GM]}}
                \texttt{HAMLE: P, (5, 1)} \\
            }
        }
    }
    & & \\ \\

    \theutterance \stepcounter{utterance}  
    & & & \multicolumn{2}{p{0.3\linewidth}}{
        \cellcolor[rgb]{0.9,0.9,0.9}{
            \makecell[{{p{\linewidth}}}]{
                \texttt{\tiny{[GM$|$GM]}}
                \texttt{Ceza: (5,1) boş değil; "│" içeriyor.} \\
\\ 
\texttt{Toplamda 16 cezanın 12 kadarını biriktirdiniz.Lütfen tekrar deneyin!} \\
            }
        }
    }
    & & \\ \\

    \theutterance \stepcounter{utterance}  
    & & \multicolumn{4}{p{0.6\linewidth}}{
        \cellcolor[rgb]{0.9,0.9,0.9}{
            \makecell[{{p{\linewidth}}}]{
                \texttt{\tiny{[P1$\langle$GM]}}
                \texttt{Ceza: (5,1) boş değil; "│" içeriyor.} \\
\\ 
\texttt{Toplamda 16 cezanın 12 kadarını biriktirdiniz.Lütfen tekrar deneyin!} \\
            }
        }
    }
    & \\ \\

    \theutterance \stepcounter{utterance}  
    & \multicolumn{4}{p{0.6\linewidth}}{
        \cellcolor[rgb]{0.9,0.9,0.9}{
            \makecell[{{p{\linewidth}}}]{
                \texttt{\tiny{[P1$\rangle$GM]}}
                \texttt{HAMLE: P, (6, 1)} \\
            }
        }
    }
    & & \\ \\

    \theutterance \stepcounter{utterance}  
    & & & \multicolumn{2}{p{0.3\linewidth}}{
        \cellcolor[rgb]{0.9,0.9,0.9}{
            \makecell[{{p{\linewidth}}}]{
                \texttt{\tiny{[GM$|$GM]}}
                \texttt{HAMLE: P, (6, 1)} \\
            }
        }
    }
    & & \\ \\

    \theutterance \stepcounter{utterance}  
    & & & \multicolumn{2}{p{0.3\linewidth}}{
        \cellcolor[rgb]{0.9,0.9,0.9}{
            \makecell[{{p{\linewidth}}}]{
                \texttt{\tiny{[GM$|$GM]}}
                \texttt{Ceza: (6,1) boş değil; "P" içeriyor.} \\
\\ 
\texttt{Toplamda 16 cezanın 13 kadarını biriktirdiniz.Lütfen tekrar deneyin!} \\
            }
        }
    }
    & & \\ \\

    \theutterance \stepcounter{utterance}  
    & & \multicolumn{4}{p{0.6\linewidth}}{
        \cellcolor[rgb]{0.9,0.9,0.9}{
            \makecell[{{p{\linewidth}}}]{
                \texttt{\tiny{[P1$\langle$GM]}}
                \texttt{Ceza: (6,1) boş değil; "P" içeriyor.} \\
\\ 
\texttt{Toplamda 16 cezanın 13 kadarını biriktirdiniz.Lütfen tekrar deneyin!} \\
            }
        }
    }
    & \\ \\

    \theutterance \stepcounter{utterance}  
    & \multicolumn{4}{p{0.6\linewidth}}{
        \cellcolor[rgb]{0.9,0.9,0.9}{
            \makecell[{{p{\linewidth}}}]{
                \texttt{\tiny{[P1$\rangle$GM]}}
                \texttt{HAMLE: P, (5, 2)} \\
            }
        }
    }
    & & \\ \\

    \theutterance \stepcounter{utterance}  
    & & & \multicolumn{2}{p{0.3\linewidth}}{
        \cellcolor[rgb]{0.9,0.9,0.9}{
            \makecell[{{p{\linewidth}}}]{
                \texttt{\tiny{[GM$|$GM]}}
                \texttt{HAMLE: P, (5, 2)} \\
            }
        }
    }
    & & \\ \\

    \theutterance \stepcounter{utterance}  
    & & & \multicolumn{2}{p{0.3\linewidth}}{
        \cellcolor[rgb]{0.9,0.9,0.9}{
            \makecell[{{p{\linewidth}}}]{
                \texttt{\tiny{[GM$|$GM]}}
                \texttt{Ceza: (5,2) boş değil; "├" içeriyor.} \\
\\ 
\texttt{Toplamda 16 cezanın 14 kadarını biriktirdiniz.Lütfen tekrar deneyin!} \\
            }
        }
    }
    & & \\ \\

    \theutterance \stepcounter{utterance}  
    & & \multicolumn{4}{p{0.6\linewidth}}{
        \cellcolor[rgb]{0.9,0.9,0.9}{
            \makecell[{{p{\linewidth}}}]{
                \texttt{\tiny{[P1$\langle$GM]}}
                \texttt{Ceza: (5,2) boş değil; "├" içeriyor.} \\
\\ 
\texttt{Toplamda 16 cezanın 14 kadarını biriktirdiniz.Lütfen tekrar deneyin!} \\
            }
        }
    }
    & \\ \\

    \theutterance \stepcounter{utterance}  
    & \multicolumn{4}{p{0.6\linewidth}}{
        \cellcolor[rgb]{0.9,0.9,0.9}{
            \makecell[{{p{\linewidth}}}]{
                \texttt{\tiny{[P1$\rangle$GM]}}
                \texttt{HAMLE: P, (4, 2)} \\
            }
        }
    }
    & & \\ \\

    \theutterance \stepcounter{utterance}  
    & & & \multicolumn{2}{p{0.3\linewidth}}{
        \cellcolor[rgb]{0.9,0.9,0.9}{
            \makecell[{{p{\linewidth}}}]{
                \texttt{\tiny{[GM$|$GM]}}
                \texttt{HAMLE: P, (4, 2)} \\
            }
        }
    }
    & & \\ \\

    \theutterance \stepcounter{utterance}  
    & & & \multicolumn{2}{p{0.3\linewidth}}{
        \cellcolor[rgb]{0.9,0.9,0.9}{
            \makecell[{{p{\linewidth}}}]{
                \texttt{\tiny{[GM$|$GM]}}
                \texttt{P moved to (4, 2)} \\
            }
        }
    }
    & & \\ \\

    \theutterance \stepcounter{utterance}  
    & & \multicolumn{4}{p{0.6\linewidth}}{
        \cellcolor[rgb]{0.9,0.9,0.9}{
            \makecell[{{p{\linewidth}}}]{
                \texttt{\tiny{[GM$\rangle$P2]}}
                \texttt{Mesajın diğer oyuncuya iletildi.} \\
\\ 
\\ 
\texttt{Şu anda en fazla 28 turdan 9. turu oynuyorsunuz.} \\
\texttt{Toplamda 16 cezanın 14 kadarını biriktirdiniz.Diğer oyuncu kendi ızgarasında bir nesneyi hareket ettirdi.} \\
\\ 
\\ 
\texttt{Bir sonraki komutun nedir?} \\
            }
        }
    }
    & \\ \\

    \theutterance \stepcounter{utterance}  
    & & & \multicolumn{4}{p{0.6\linewidth}}{
        \cellcolor[rgb]{0.9,0.9,0.9}{
            \makecell[{{p{\linewidth}}}]{
                \texttt{\tiny{[GM$\langle$P2]}}
                \texttt{HAMLE: P, (2,1)} \\
            }
        }
    }
     \\ \\

    \theutterance \stepcounter{utterance}  
    & & & \multicolumn{2}{p{0.3\linewidth}}{
        \cellcolor[rgb]{0.9,0.9,0.9}{
            \makecell[{{p{\linewidth}}}]{
                \texttt{\tiny{[GM$|$GM]}}
                \texttt{HAMLE: P, (2,1)} \\
            }
        }
    }
    & & \\ \\

    \theutterance \stepcounter{utterance}  
    & & & \multicolumn{2}{p{0.3\linewidth}}{
        \cellcolor[rgb]{0.9,0.9,0.9}{
            \makecell[{{p{\linewidth}}}]{
                \texttt{\tiny{[GM$|$GM]}}
                \texttt{P moved to (2, 1)} \\
            }
        }
    }
    & & \\ \\

    \theutterance \stepcounter{utterance}  
    & & \multicolumn{4}{p{0.6\linewidth}}{
        \cellcolor[rgb]{0.9,0.9,0.9}{
            \makecell[{{p{\linewidth}}}]{
                \texttt{\tiny{[P1$\langle$GM]}}
                \texttt{P’ başarıyla (4,2) koordinatlarına taşındı. Güncellenmiş ızgaran şu şekilde:} \\
\texttt{```} \\
\texttt{ 1234567} \\
\texttt{╔════╤══╗} \\
\texttt{║◌◌◌O│◌M║ 1} \\
\texttt{║◌◌TP├──╢ 2} \\
\texttt{║S◌◌◌│◌◌║ 3} \\
\texttt{║◌┌──┼──╢ 4} \\
\texttt{║◌│◌A│D◌║ 5} \\
\texttt{╟─┴┬─┤◌◌║ 6} \\
\texttt{║◌◌│◌│◌◌║ 7} \\
\texttt{╚══╧═╧══╝} \\
\texttt{```} \\
\texttt{Şu anda en fazla 28 turdan 9. turu oynuyorsunuz.} \\
\texttt{Toplamda 16 cezanın 14 kadarını biriktirdiniz.Diğer oyuncu kendi ızgarasında bir nesneyi hareket ettirdi.} \\
\\ 
\\ 
\texttt{Bir sonraki komutun nedir?} \\
            }
        }
    }
    & \\ \\

    \theutterance \stepcounter{utterance}  
    & \multicolumn{4}{p{0.6\linewidth}}{
        \cellcolor[rgb]{0.9,0.9,0.9}{
            \makecell[{{p{\linewidth}}}]{
                \texttt{\tiny{[P1$\rangle$GM]}}
                \texttt{SÖYLE: P nesnesini (4,2) konumuna başarıyla taşıdım. Ceza sınırına çok yaklaştığımız için çok dikkatli olmalıyım. Diğer oyuncu hangi nesneyi hareket ettirdi? Bir sonraki hamlemde M nesnesini merkeze yaklaştırmayı planlıyorum.} \\
            }
        }
    }
    & & \\ \\

    \theutterance \stepcounter{utterance}  
    & & & \multicolumn{2}{p{0.3\linewidth}}{
        \cellcolor[rgb]{0.9,0.9,0.9}{
            \makecell[{{p{\linewidth}}}]{
                \texttt{\tiny{[GM$|$GM]}}
                \texttt{SÖYLE: P nesnesini (4,2) konumuna başarıyla taşıdım. Ceza sınırına çok yaklaştığımız için çok dikkatli olmalıyım. Diğer oyuncu hangi nesneyi hareket ettirdi? Bir sonraki hamlemde M nesnesini merkeze yaklaştırmayı planlıyorum.} \\
            }
        }
    }
    & & \\ \\

    \theutterance \stepcounter{utterance}  
    & & \multicolumn{4}{p{0.6\linewidth}}{
        \cellcolor[rgb]{0.9,0.9,0.9}{
            \makecell[{{p{\linewidth}}}]{
                \texttt{\tiny{[GM$\rangle$P2]}}
                \texttt{P’ başarıyla (2,1) koordinatlarına taşındı. Güncellenmiş ızgaran şu şekilde:} \\
\texttt{```} \\
\texttt{ 1234567} \\
\texttt{╔════╤══╗} \\
\texttt{║◌P◌◌│O◌║ 1} \\
\texttt{║◌SM◌├──╢ 2} \\
\texttt{║◌◌◌◌│T◌║ 3} \\
\texttt{║A┌──┼──╢ 4} \\
\texttt{║◌│◌◌│◌◌║ 5} \\
\texttt{╟─┴┬─┤◌◌║ 6} \\
\texttt{║◌◌│◌│D◌║ 7} \\
\texttt{╚══╧═╧══╝} \\
\texttt{```} \\
\texttt{Şu anda en fazla 28 turdan 10. turu oynuyorsunuz.} \\
\texttt{Toplamda 16 cezanın 14 kadarını biriktirdiniz.Diğer oyuncu şu mesajı gönderdi:} \\
\\ 
\\ 
\texttt{"P nesnesini (4,2) konumuna başarıyla taşıdım. Ceza sınırına çok yaklaştığımız için çok dikkatli olmalıyım. Diğer oyuncu hangi nesneyi hareket ettirdi? Bir sonraki hamlemde M nesnesini merkeze yaklaştırmayı planlıyorum."} \\
\\ 
\\ 
\texttt{Bir sonraki komutun nedir?} \\
            }
        }
    }
    & \\ \\

    \theutterance \stepcounter{utterance}  
    & & & \multicolumn{4}{p{0.6\linewidth}}{
        \cellcolor[rgb]{0.9,0.9,0.9}{
            \makecell[{{p{\linewidth}}}]{
                \texttt{\tiny{[GM$\langle$P2]}}
                \texttt{SÖYLE: Ben P nesnesini (2,1) konumuna taşıdım. Haklısın, ceza sınırına çok yaklaştık, çok dikkatli olmalıyız. M nesnesini merkeze yaklaştırma planına devam edebilirsin. Ben de bir sonraki hamlemde M nesnesini hareket ettireceğim.} \\
            }
        }
    }
     \\ \\

    \theutterance \stepcounter{utterance}  
    & & & \multicolumn{2}{p{0.3\linewidth}}{
        \cellcolor[rgb]{0.9,0.9,0.9}{
            \makecell[{{p{\linewidth}}}]{
                \texttt{\tiny{[GM$|$GM]}}
                \texttt{SÖYLE: Ben P nesnesini (2,1) konumuna taşıdım. Haklısın, ceza sınırına çok yaklaştık, çok dikkatli olmalıyız. M nesnesini merkeze yaklaştırma planına devam edebilirsin. Ben de bir sonraki hamlemde M nesnesini hareket ettireceğim.} \\
            }
        }
    }
    & & \\ \\

    \theutterance \stepcounter{utterance}  
    & & \multicolumn{4}{p{0.6\linewidth}}{
        \cellcolor[rgb]{0.9,0.9,0.9}{
            \makecell[{{p{\linewidth}}}]{
                \texttt{\tiny{[P1$\langle$GM]}}
                \texttt{Mesajın diğer oyuncuya iletildi.} \\
\\ 
\\ 
\texttt{Şu anda en fazla 28 turdan 10. turu oynuyorsunuz.} \\
\texttt{Toplamda 16 cezanın 14 kadarını biriktirdiniz.Diğer oyuncu şu mesajı gönderdi:} \\
\\ 
\\ 
\texttt{"Ben P nesnesini (2,1) konumuna taşıdım. Haklısın, ceza sınırına çok yaklaştık, çok dikkatli olmalıyız. M nesnesini merkeze yaklaştırma planına devam edebilirsin. Ben de bir sonraki hamlemde M nesnesini hareket ettireceğim."} \\
\\ 
\\ 
\texttt{Bir sonraki komutun nedir?} \\
            }
        }
    }
    & \\ \\

    \theutterance \stepcounter{utterance}  
    & \multicolumn{4}{p{0.6\linewidth}}{
        \cellcolor[rgb]{0.9,0.9,0.9}{
            \makecell[{{p{\linewidth}}}]{
                \texttt{\tiny{[P1$\rangle$GM]}}
                \texttt{HAMLE: M, (5, 1)} \\
            }
        }
    }
    & & \\ \\

    \theutterance \stepcounter{utterance}  
    & & & \multicolumn{2}{p{0.3\linewidth}}{
        \cellcolor[rgb]{0.9,0.9,0.9}{
            \makecell[{{p{\linewidth}}}]{
                \texttt{\tiny{[GM$|$GM]}}
                \texttt{HAMLE: M, (5, 1)} \\
            }
        }
    }
    & & \\ \\

    \theutterance \stepcounter{utterance}  
    & & & \multicolumn{2}{p{0.3\linewidth}}{
        \cellcolor[rgb]{0.9,0.9,0.9}{
            \makecell[{{p{\linewidth}}}]{
                \texttt{\tiny{[GM$|$GM]}}
                \texttt{Ceza: (5,1) boş değil; "│" içeriyor.} \\
\\ 
\texttt{Toplamda 16 cezanın 15 kadarını biriktirdiniz.Lütfen tekrar deneyin!} \\
            }
        }
    }
    & & \\ \\

    \theutterance \stepcounter{utterance}  
    & & \multicolumn{4}{p{0.6\linewidth}}{
        \cellcolor[rgb]{0.9,0.9,0.9}{
            \makecell[{{p{\linewidth}}}]{
                \texttt{\tiny{[P1$\langle$GM]}}
                \texttt{Ceza: (5,1) boş değil; "│" içeriyor.} \\
\\ 
\texttt{Toplamda 16 cezanın 15 kadarını biriktirdiniz.Lütfen tekrar deneyin!} \\
            }
        }
    }
    & \\ \\

    \theutterance \stepcounter{utterance}  
    & \multicolumn{4}{p{0.6\linewidth}}{
        \cellcolor[rgb]{0.9,0.9,0.9}{
            \makecell[{{p{\linewidth}}}]{
                \texttt{\tiny{[P1$\rangle$GM]}}
                \texttt{HAMLE: M, (6, 1)} \\
            }
        }
    }
    & & \\ \\

    \theutterance \stepcounter{utterance}  
    & & & \multicolumn{2}{p{0.3\linewidth}}{
        \cellcolor[rgb]{0.9,0.9,0.9}{
            \makecell[{{p{\linewidth}}}]{
                \texttt{\tiny{[GM$|$GM]}}
                \texttt{HAMLE: M, (6, 1)} \\
            }
        }
    }
    & & \\ \\

    \theutterance \stepcounter{utterance}  
    & & & \multicolumn{2}{p{0.3\linewidth}}{
        \cellcolor[rgb]{0.9,0.9,0.9}{
            \makecell[{{p{\linewidth}}}]{
                \texttt{\tiny{[GM$|$GM]}}
                \texttt{M moved to (6, 1)} \\
            }
        }
    }
    & & \\ \\

    \theutterance \stepcounter{utterance}  
    & & \multicolumn{4}{p{0.6\linewidth}}{
        \cellcolor[rgb]{0.9,0.9,0.9}{
            \makecell[{{p{\linewidth}}}]{
                \texttt{\tiny{[GM$\rangle$P2]}}
                \texttt{Mesajın diğer oyuncuya iletildi.} \\
\\ 
\\ 
\texttt{Şu anda en fazla 28 turdan 11. turu oynuyorsunuz.} \\
\texttt{Toplamda 16 cezanın 15 kadarını biriktirdiniz.Diğer oyuncu kendi ızgarasında bir nesneyi hareket ettirdi.} \\
\\ 
\\ 
\texttt{Bir sonraki komutun nedir?} \\
            }
        }
    }
    & \\ \\

    \theutterance \stepcounter{utterance}  
    & & & \multicolumn{4}{p{0.6\linewidth}}{
        \cellcolor[rgb]{0.9,0.9,0.9}{
            \makecell[{{p{\linewidth}}}]{
                \texttt{\tiny{[GM$\langle$P2]}}
                \texttt{HAMLE: M, (2,2)} \\
            }
        }
    }
     \\ \\

    \theutterance \stepcounter{utterance}  
    & & & \multicolumn{2}{p{0.3\linewidth}}{
        \cellcolor[rgb]{0.9,0.9,0.9}{
            \makecell[{{p{\linewidth}}}]{
                \texttt{\tiny{[GM$|$GM]}}
                \texttt{HAMLE: M, (2,2)} \\
            }
        }
    }
    & & \\ \\

    \theutterance \stepcounter{utterance}  
    & & & \multicolumn{2}{p{0.3\linewidth}}{
        \cellcolor[rgb]{0.9,0.9,0.9}{
            \makecell[{{p{\linewidth}}}]{
                \texttt{\tiny{[GM$|$GM]}}
                \texttt{Ceza: (2,2) boş değil; "S" içeriyor.} \\
\\ 
\texttt{Toplamda 16 cezanın 16 kadarını biriktirdiniz.Lütfen tekrar deneyin!} \\
            }
        }
    }
    & & \\ \\

    \theutterance \stepcounter{utterance}  
    & & \multicolumn{4}{p{0.6\linewidth}}{
        \cellcolor[rgb]{0.9,0.9,0.9}{
            \makecell[{{p{\linewidth}}}]{
                \texttt{\tiny{[GM$\rangle$P2]}}
                \texttt{Ceza: (2,2) boş değil; "S" içeriyor.} \\
\\ 
\texttt{Toplamda 16 cezanın 16 kadarını biriktirdiniz.Lütfen tekrar deneyin!} \\
            }
        }
    }
    & \\ \\

    \theutterance \stepcounter{utterance}  
    & & & \multicolumn{4}{p{0.6\linewidth}}{
        \cellcolor[rgb]{0.9,0.9,0.9}{
            \makecell[{{p{\linewidth}}}]{
                \texttt{\tiny{[GM$\langle$P2]}}
                \texttt{HAMLE: M, (3,1)} \\
            }
        }
    }
     \\ \\

    \theutterance \stepcounter{utterance}  
    & & & \multicolumn{2}{p{0.3\linewidth}}{
        \cellcolor[rgb]{0.9,0.9,0.9}{
            \makecell[{{p{\linewidth}}}]{
                \texttt{\tiny{[GM$|$GM]}}
                \texttt{HAMLE: M, (3,1)} \\
            }
        }
    }
    & & \\ \\

    \theutterance \stepcounter{utterance}  
    & & & \multicolumn{2}{p{0.3\linewidth}}{
        \cellcolor[rgb]{0.9,0.9,0.9}{
            \makecell[{{p{\linewidth}}}]{
                \texttt{\tiny{[GM$|$GM]}}
                \texttt{M moved to (3, 1)} \\
            }
        }
    }
    & & \\ \\

    \theutterance \stepcounter{utterance}  
    & & \multicolumn{4}{p{0.6\linewidth}}{
        \cellcolor[rgb]{0.9,0.9,0.9}{
            \makecell[{{p{\linewidth}}}]{
                \texttt{\tiny{[P1$\langle$GM]}}
                \texttt{M’ başarıyla (6,1) koordinatlarına taşındı. Güncellenmiş ızgaran şu şekilde:} \\
\texttt{```} \\
\texttt{ 1234567} \\
\texttt{╔════╤══╗} \\
\texttt{║◌◌◌O│M◌║ 1} \\
\texttt{║◌◌TP├──╢ 2} \\
\texttt{║S◌◌◌│◌◌║ 3} \\
\texttt{║◌┌──┼──╢ 4} \\
\texttt{║◌│◌A│D◌║ 5} \\
\texttt{╟─┴┬─┤◌◌║ 6} \\
\texttt{║◌◌│◌│◌◌║ 7} \\
\texttt{╚══╧═╧══╝} \\
\texttt{```} \\
\texttt{Şu anda en fazla 28 turdan 11. turu oynuyorsunuz.} \\
\texttt{Toplamda 16 cezanın 16 kadarını biriktirdiniz.Diğer oyuncu kendi ızgarasında bir nesneyi hareket ettirdi.} \\
\\ 
\\ 
\texttt{Bir sonraki komutun nedir?} \\
            }
        }
    }
    & \\ \\

    \theutterance \stepcounter{utterance}  
    & \multicolumn{4}{p{0.6\linewidth}}{
        \cellcolor[rgb]{0.9,0.9,0.9}{
            \makecell[{{p{\linewidth}}}]{
                \texttt{\tiny{[P1$\rangle$GM]}}
                \texttt{SÖYLE: M nesnesini (6,1) konumuna taşıdım. Çok kritik bir noktadayız, tek bir hata daha oyunu kaybetmemize neden olur. Lütfen çok dikkatli olalım. Diğer oyuncu hangi nesneyi hareket ettirdi?} \\
            }
        }
    }
    & & \\ \\

    \theutterance \stepcounter{utterance}  
    & & & \multicolumn{2}{p{0.3\linewidth}}{
        \cellcolor[rgb]{0.9,0.9,0.9}{
            \makecell[{{p{\linewidth}}}]{
                \texttt{\tiny{[GM$|$GM]}}
                \texttt{SÖYLE: M nesnesini (6,1) konumuna taşıdım. Çok kritik bir noktadayız, tek bir hata daha oyunu kaybetmemize neden olur. Lütfen çok dikkatli olalım. Diğer oyuncu hangi nesneyi hareket ettirdi?} \\
            }
        }
    }
    & & \\ \\

    \theutterance \stepcounter{utterance}  
    & & \multicolumn{4}{p{0.6\linewidth}}{
        \cellcolor[rgb]{0.9,0.9,0.9}{
            \makecell[{{p{\linewidth}}}]{
                \texttt{\tiny{[GM$\rangle$P2]}}
                \texttt{M’ başarıyla (3,1) koordinatlarına taşındı. Güncellenmiş ızgaran şu şekilde:} \\
\texttt{```} \\
\texttt{ 1234567} \\
\texttt{╔════╤══╗} \\
\texttt{║◌PM◌│O◌║ 1} \\
\texttt{║◌S◌◌├──╢ 2} \\
\texttt{║◌◌◌◌│T◌║ 3} \\
\texttt{║A┌──┼──╢ 4} \\
\texttt{║◌│◌◌│◌◌║ 5} \\
\texttt{╟─┴┬─┤◌◌║ 6} \\
\texttt{║◌◌│◌│D◌║ 7} \\
\texttt{╚══╧═╧══╝} \\
\texttt{```} \\
\texttt{Şu anda en fazla 28 turdan 12. turu oynuyorsunuz.} \\
\texttt{Toplamda 16 cezanın 16 kadarını biriktirdiniz.Diğer oyuncu şu mesajı gönderdi:} \\
\\ 
\\ 
\texttt{"M nesnesini (6,1) konumuna taşıdım. Çok kritik bir noktadayız, tek bir hata daha oyunu kaybetmemize neden olur. Lütfen çok dikkatli olalım. Diğer oyuncu hangi nesneyi hareket ettirdi?"} \\
\\ 
\\ 
\texttt{Bir sonraki komutun nedir?} \\
            }
        }
    }
    & \\ \\

    \theutterance \stepcounter{utterance}  
    & & & \multicolumn{4}{p{0.6\linewidth}}{
        \cellcolor[rgb]{0.9,0.9,0.9}{
            \makecell[{{p{\linewidth}}}]{
                \texttt{\tiny{[GM$\langle$P2]}}
                \texttt{SÖYLE: Ben M nesnesini (3,1) konumuna taşıdım. Haklısın, durum çok kritik. Tek bir hata bile bizi kaybettirebilir. Bundan sonraki tüm hamlelerimi yapmadan önce hücrenin tamamen boş olduğundan emin olmak için çok dikkatli olacağım. Bir sonraki hamlemde O nesnesini hareket ettirmeyi planlıyorum.} \\
            }
        }
    }
     \\ \\

    \theutterance \stepcounter{utterance}  
    & & & \multicolumn{2}{p{0.3\linewidth}}{
        \cellcolor[rgb]{0.9,0.9,0.9}{
            \makecell[{{p{\linewidth}}}]{
                \texttt{\tiny{[GM$|$GM]}}
                \texttt{SÖYLE: Ben M nesnesini (3,1) konumuna taşıdım. Haklısın, durum çok kritik. Tek bir hata bile bizi kaybettirebilir. Bundan sonraki tüm hamlelerimi yapmadan önce hücrenin tamamen boş olduğundan emin olmak için çok dikkatli olacağım. Bir sonraki hamlemde O nesnesini hareket ettirmeyi planlıyorum.} \\
            }
        }
    }
    & & \\ \\

    \theutterance \stepcounter{utterance}  
    & & \multicolumn{4}{p{0.6\linewidth}}{
        \cellcolor[rgb]{0.9,0.9,0.9}{
            \makecell[{{p{\linewidth}}}]{
                \texttt{\tiny{[P1$\langle$GM]}}
                \texttt{Mesajın diğer oyuncuya iletildi.} \\
\\ 
\\ 
\texttt{Şu anda en fazla 28 turdan 12. turu oynuyorsunuz.} \\
\texttt{Toplamda 16 cezanın 16 kadarını biriktirdiniz.Diğer oyuncu şu mesajı gönderdi:} \\
\\ 
\\ 
\texttt{"Ben M nesnesini (3,1) konumuna taşıdım. Haklısın, durum çok kritik. Tek bir hata bile bizi kaybettirebilir. Bundan sonraki tüm hamlelerimi yapmadan önce hücrenin tamamen boş olduğundan emin olmak için çok dikkatli olacağım. Bir sonraki hamlemde O nesnesini hareket ettirmeyi planlıyorum."} \\
\\ 
\\ 
\texttt{Bir sonraki komutun nedir?} \\
            }
        }
    }
    & \\ \\

    \theutterance \stepcounter{utterance}  
    & \multicolumn{4}{p{0.6\linewidth}}{
        \cellcolor[rgb]{0.9,0.9,0.9}{
            \makecell[{{p{\linewidth}}}]{
                \texttt{\tiny{[P1$\rangle$GM]}}
                \texttt{SÖYLE: M nesnesini (6,1) konumuna taşıdım. Çok dikkatli olmalıyız, hata payımız kalmadı. O nesnesini hareket ettirme planını bekliyorum. Ben de bir sonraki hamlemde T nesnesini merkeze yaklaştırmayı düşüneceğim.} \\
            }
        }
    }
    & & \\ \\

    \theutterance \stepcounter{utterance}  
    & & & \multicolumn{2}{p{0.3\linewidth}}{
        \cellcolor[rgb]{0.9,0.9,0.9}{
            \makecell[{{p{\linewidth}}}]{
                \texttt{\tiny{[GM$|$GM]}}
                \texttt{SÖYLE: M nesnesini (6,1) konumuna taşıdım. Çok dikkatli olmalıyız, hata payımız kalmadı. O nesnesini hareket ettirme planını bekliyorum. Ben de bir sonraki hamlemde T nesnesini merkeze yaklaştırmayı düşüneceğim.} \\
            }
        }
    }
    & & \\ \\

    \theutterance \stepcounter{utterance}  
    & & \multicolumn{4}{p{0.6\linewidth}}{
        \cellcolor[rgb]{0.9,0.9,0.9}{
            \makecell[{{p{\linewidth}}}]{
                \texttt{\tiny{[GM$\rangle$P2]}}
                \texttt{Mesajın diğer oyuncuya iletildi.} \\
\\ 
\\ 
\texttt{Şu anda en fazla 28 turdan 13. turu oynuyorsunuz.} \\
\texttt{Toplamda 16 cezanın 16 kadarını biriktirdiniz.Diğer oyuncu şu mesajı gönderdi:} \\
\\ 
\\ 
\texttt{"M nesnesini (6,1) konumuna taşıdım. Çok dikkatli olmalıyız, hata payımız kalmadı. O nesnesini hareket ettirme planını bekliyorum. Ben de bir sonraki hamlemde T nesnesini merkeze yaklaştırmayı düşüneceğim."} \\
\\ 
\\ 
\texttt{Bir sonraki komutun nedir?} \\
            }
        }
    }
    & \\ \\

    \theutterance \stepcounter{utterance}  
    & & & \multicolumn{4}{p{0.6\linewidth}}{
        \cellcolor[rgb]{0.9,0.9,0.9}{
            \makecell[{{p{\linewidth}}}]{
                \texttt{\tiny{[GM$\langle$P2]}}
                \texttt{HAMLE: O, (4,1)} \\
            }
        }
    }
     \\ \\

    \theutterance \stepcounter{utterance}  
    & & & \multicolumn{2}{p{0.3\linewidth}}{
        \cellcolor[rgb]{0.9,0.9,0.9}{
            \makecell[{{p{\linewidth}}}]{
                \texttt{\tiny{[GM$|$GM]}}
                \texttt{HAMLE: O, (4,1)} \\
            }
        }
    }
    & & \\ \\

    \theutterance \stepcounter{utterance}  
    & & & \multicolumn{2}{p{0.3\linewidth}}{
        \cellcolor[rgb]{0.9,0.9,0.9}{
            \makecell[{{p{\linewidth}}}]{
                \texttt{\tiny{[GM$|$GM]}}
                \texttt{O moved to (4, 1)} \\
            }
        }
    }
    & & \\ \\

    \theutterance \stepcounter{utterance}  
    & & \multicolumn{4}{p{0.6\linewidth}}{
        \cellcolor[rgb]{0.9,0.9,0.9}{
            \makecell[{{p{\linewidth}}}]{
                \texttt{\tiny{[P1$\langle$GM]}}
                \texttt{Mesajın diğer oyuncuya iletildi.} \\
\\ 
\\ 
\texttt{Şu anda en fazla 28 turdan 13. turu oynuyorsunuz.} \\
\texttt{Toplamda 16 cezanın 16 kadarını biriktirdiniz.Diğer oyuncu kendi ızgarasında bir nesneyi hareket ettirdi.} \\
\\ 
\\ 
\texttt{Bir sonraki komutun nedir?} \\
            }
        }
    }
    & \\ \\

    \theutterance \stepcounter{utterance}  
    & \multicolumn{4}{p{0.6\linewidth}}{
        \cellcolor[rgb]{0.9,0.9,0.9}{
            \makecell[{{p{\linewidth}}}]{
                \texttt{\tiny{[P1$\rangle$GM]}}
                \texttt{SÖYLE: Diğer oyuncu O nesnesini hareket ettirdi mi? Lütfen hangi nesneyi nereye taşıdığını söyler misin? Ben de bir sonraki hamlemde T nesnesini hareket ettireceğim ama hata yapmamak için tahtayı çok dikkatli incelemem gerekiyor.} \\
            }
        }
    }
    & & \\ \\

    \theutterance \stepcounter{utterance}  
    & & & \multicolumn{2}{p{0.3\linewidth}}{
        \cellcolor[rgb]{0.9,0.9,0.9}{
            \makecell[{{p{\linewidth}}}]{
                \texttt{\tiny{[GM$|$GM]}}
                \texttt{SÖYLE: Diğer oyuncu O nesnesini hareket ettirdi mi? Lütfen hangi nesneyi nereye taşıdığını söyler misin? Ben de bir sonraki hamlemde T nesnesini hareket ettireceğim ama hata yapmamak için tahtayı çok dikkatli incelemem gerekiyor.} \\
            }
        }
    }
    & & \\ \\

    \theutterance \stepcounter{utterance}  
    & & \multicolumn{4}{p{0.6\linewidth}}{
        \cellcolor[rgb]{0.9,0.9,0.9}{
            \makecell[{{p{\linewidth}}}]{
                \texttt{\tiny{[GM$\rangle$P2]}}
                \texttt{O’ başarıyla (4,1) koordinatlarına taşındı. Güncellenmiş ızgaran şu şekilde:} \\
\texttt{```} \\
\texttt{ 1234567} \\
\texttt{╔════╤══╗} \\
\texttt{║◌PMO│◌◌║ 1} \\
\texttt{║◌S◌◌├──╢ 2} \\
\texttt{║◌◌◌◌│T◌║ 3} \\
\texttt{║A┌──┼──╢ 4} \\
\texttt{║◌│◌◌│◌◌║ 5} \\
\texttt{╟─┴┬─┤◌◌║ 6} \\
\texttt{║◌◌│◌│D◌║ 7} \\
\texttt{╚══╧═╧══╝} \\
\texttt{```} \\
\texttt{Şu anda en fazla 28 turdan 14. turu oynuyorsunuz.} \\
\texttt{Toplamda 16 cezanın 16 kadarını biriktirdiniz.Diğer oyuncu şu mesajı gönderdi:} \\
\\ 
\\ 
\texttt{"Diğer oyuncu O nesnesini hareket ettirdi mi? Lütfen hangi nesneyi nereye taşıdığını söyler misin? Ben de bir sonraki hamlemde T nesnesini hareket ettireceğim ama hata yapmamak için tahtayı çok dikkatli incelemem gerekiyor."} \\
\\ 
\\ 
\texttt{Bir sonraki komutun nedir?} \\
            }
        }
    }
    & \\ \\

    \theutterance \stepcounter{utterance}  
    & & & \multicolumn{4}{p{0.6\linewidth}}{
        \cellcolor[rgb]{0.9,0.9,0.9}{
            \makecell[{{p{\linewidth}}}]{
                \texttt{\tiny{[GM$\langle$P2]}}
                \texttt{SÖYLE: Evet, O nesnesini (4,1) konumuna taşıdım. Çok dikkatli olalım, hata yapmamak için her hamleden önce hücreyi tekrar kontrol edeceğim. Sen T nesnesini hareket ettirdiğinde ben de bir sonraki hamlem için planımı netleştireceğim.} \\
            }
        }
    }
     \\ \\

    \theutterance \stepcounter{utterance}  
    & & & \multicolumn{2}{p{0.3\linewidth}}{
        \cellcolor[rgb]{0.9,0.9,0.9}{
            \makecell[{{p{\linewidth}}}]{
                \texttt{\tiny{[GM$|$GM]}}
                \texttt{SÖYLE: Evet, O nesnesini (4,1) konumuna taşıdım. Çok dikkatli olalım, hata yapmamak için her hamleden önce hücreyi tekrar kontrol edeceğim. Sen T nesnesini hareket ettirdiğinde ben de bir sonraki hamlem için planımı netleştireceğim.} \\
            }
        }
    }
    & & \\ \\

    \theutterance \stepcounter{utterance}  
    & & \multicolumn{4}{p{0.6\linewidth}}{
        \cellcolor[rgb]{0.9,0.9,0.9}{
            \makecell[{{p{\linewidth}}}]{
                \texttt{\tiny{[P1$\langle$GM]}}
                \texttt{Mesajın diğer oyuncuya iletildi.} \\
\\ 
\\ 
\texttt{Şu anda en fazla 28 turdan 14. turu oynuyorsunuz.} \\
\texttt{Toplamda 16 cezanın 16 kadarını biriktirdiniz.Diğer oyuncu şu mesajı gönderdi:} \\
\\ 
\\ 
\texttt{"Evet, O nesnesini (4,1) konumuna taşıdım. Çok dikkatli olalım, hata yapmamak için her hamleden önce hücreyi tekrar kontrol edeceğim. Sen T nesnesini hareket ettirdiğinde ben de bir sonraki hamlem için planımı netleştireceğim."} \\
\\ 
\\ 
\texttt{Bir sonraki komutun nedir?} \\
            }
        }
    }
    & \\ \\

    \theutterance \stepcounter{utterance}  
    & \multicolumn{4}{p{0.6\linewidth}}{
        \cellcolor[rgb]{0.9,0.9,0.9}{
            \makecell[{{p{\linewidth}}}]{
                \texttt{\tiny{[P1$\rangle$GM]}}
                \texttt{HAMLE: T, (3, 3)} \\
            }
        }
    }
    & & \\ \\

    \theutterance \stepcounter{utterance}  
    & & & \multicolumn{2}{p{0.3\linewidth}}{
        \cellcolor[rgb]{0.9,0.9,0.9}{
            \makecell[{{p{\linewidth}}}]{
                \texttt{\tiny{[GM$|$GM]}}
                \texttt{HAMLE: T, (3, 3)} \\
            }
        }
    }
    & & \\ \\

    \theutterance \stepcounter{utterance}  
    & & & \multicolumn{2}{p{0.3\linewidth}}{
        \cellcolor[rgb]{0.9,0.9,0.9}{
            \makecell[{{p{\linewidth}}}]{
                \texttt{\tiny{[GM$|$GM]}}
                \texttt{T moved to (3, 3)} \\
            }
        }
    }
    & & \\ \\

    \theutterance \stepcounter{utterance}  
    & & \multicolumn{4}{p{0.6\linewidth}}{
        \cellcolor[rgb]{0.9,0.9,0.9}{
            \makecell[{{p{\linewidth}}}]{
                \texttt{\tiny{[GM$\rangle$P2]}}
                \texttt{Mesajın diğer oyuncuya iletildi.} \\
\\ 
\\ 
\texttt{Şu anda en fazla 28 turdan 15. turu oynuyorsunuz.} \\
\texttt{Toplamda 16 cezanın 16 kadarını biriktirdiniz.Diğer oyuncu kendi ızgarasında bir nesneyi hareket ettirdi.} \\
\\ 
\\ 
\texttt{Bir sonraki komutun nedir?} \\
            }
        }
    }
    & \\ \\

    \theutterance \stepcounter{utterance}  
    & & & \multicolumn{4}{p{0.6\linewidth}}{
        \cellcolor[rgb]{0.9,0.9,0.9}{
            \makecell[{{p{\linewidth}}}]{
                \texttt{\tiny{[GM$\langle$P2]}}
                \texttt{SÖYLE: Ben O nesnesini (4,1) konumuna taşıdım. Sen T nesnesini nereye taşıdın? Bir sonraki hamlemde T nesnesini veya diğer nesneleri merkeze yaklaştırmayı planlıyorum. Lütfen kendi hamleni de paylaş ki koordinasyonumuzu bozmayalım.} \\
            }
        }
    }
     \\ \\

    \theutterance \stepcounter{utterance}  
    & & & \multicolumn{2}{p{0.3\linewidth}}{
        \cellcolor[rgb]{0.9,0.9,0.9}{
            \makecell[{{p{\linewidth}}}]{
                \texttt{\tiny{[GM$|$GM]}}
                \texttt{SÖYLE: Ben O nesnesini (4,1) konumuna taşıdım. Sen T nesnesini nereye taşıdın? Bir sonraki hamlemde T nesnesini veya diğer nesneleri merkeze yaklaştırmayı planlıyorum. Lütfen kendi hamleni de paylaş ki koordinasyonumuzu bozmayalım.} \\
            }
        }
    }
    & & \\ \\

    \theutterance \stepcounter{utterance}  
    & & \multicolumn{4}{p{0.6\linewidth}}{
        \cellcolor[rgb]{0.9,0.9,0.9}{
            \makecell[{{p{\linewidth}}}]{
                \texttt{\tiny{[P1$\langle$GM]}}
                \texttt{T’ başarıyla (3,3) koordinatlarına taşındı. Güncellenmiş ızgaran şu şekilde:} \\
\texttt{```} \\
\texttt{ 1234567} \\
\texttt{╔════╤══╗} \\
\texttt{║◌◌◌O│M◌║ 1} \\
\texttt{║◌◌◌P├──╢ 2} \\
\texttt{║S◌T◌│◌◌║ 3} \\
\texttt{║◌┌──┼──╢ 4} \\
\texttt{║◌│◌A│D◌║ 5} \\
\texttt{╟─┴┬─┤◌◌║ 6} \\
\texttt{║◌◌│◌│◌◌║ 7} \\
\texttt{╚══╧═╧══╝} \\
\texttt{```} \\
\texttt{Şu anda en fazla 28 turdan 15. turu oynuyorsunuz.} \\
\texttt{Toplamda 16 cezanın 16 kadarını biriktirdiniz.Diğer oyuncu şu mesajı gönderdi:} \\
\\ 
\\ 
\texttt{"Ben O nesnesini (4,1) konumuna taşıdım. Sen T nesnesini nereye taşıdın? Bir sonraki hamlemde T nesnesini veya diğer nesneleri merkeze yaklaştırmayı planlıyorum. Lütfen kendi hamleni de paylaş ki koordinasyonumuzu bozmayalım."} \\
\\ 
\\ 
\texttt{Bir sonraki komutun nedir?} \\
            }
        }
    }
    & \\ \\

    \theutterance \stepcounter{utterance}  
    & \multicolumn{4}{p{0.6\linewidth}}{
        \cellcolor[rgb]{0.9,0.9,0.9}{
            \makecell[{{p{\linewidth}}}]{
                \texttt{\tiny{[P1$\rangle$GM]}}
                \texttt{SÖYLE: T nesnesini (3,3) konumuna başarıyla taşıdım. Şu anki durumumuz şöyle: O (4,1), M (6,1), P (4,2), T (3,3), S (1,3), A (4,5), D (6,5). Bir sonraki hamlemde A nesnesini merkeze yaklaştırmak için (4,4) konumuna taşımayı deneyeceğim ama (4,4) hücresinin boş olduğundan emin olmam gerekiyor. Sen de kendi tahtandaki nesneleri bu koordinatlara göre planlayabilirsin.} \\
            }
        }
    }
    & & \\ \\

    \theutterance \stepcounter{utterance}  
    & & & \multicolumn{2}{p{0.3\linewidth}}{
        \cellcolor[rgb]{0.9,0.9,0.9}{
            \makecell[{{p{\linewidth}}}]{
                \texttt{\tiny{[GM$|$GM]}}
                \texttt{SÖYLE: T nesnesini (3,3) konumuna başarıyla taşıdım. Şu anki durumumuz şöyle: O (4,1), M (6,1), P (4,2), T (3,3), S (1,3), A (4,5), D (6,5). Bir sonraki hamlemde A nesnesini merkeze yaklaştırmak için (4,4) konumuna taşımayı deneyeceğim ama (4,4) hücresinin boş olduğundan emin olmam gerekiyor. Sen de kendi tahtandaki nesneleri bu koordinatlara göre planlayabilirsin.} \\
            }
        }
    }
    & & \\ \\

    \theutterance \stepcounter{utterance}  
    & & \multicolumn{4}{p{0.6\linewidth}}{
        \cellcolor[rgb]{0.9,0.9,0.9}{
            \makecell[{{p{\linewidth}}}]{
                \texttt{\tiny{[GM$\rangle$P2]}}
                \texttt{Mesajın diğer oyuncuya iletildi.} \\
\\ 
\\ 
\texttt{Şu anda en fazla 28 turdan 16. turu oynuyorsunuz.} \\
\texttt{Toplamda 16 cezanın 16 kadarını biriktirdiniz.Diğer oyuncu şu mesajı gönderdi:} \\
\\ 
\\ 
\texttt{"T nesnesini (3,3) konumuna başarıyla taşıdım. Şu anki durumumuz şöyle: O (4,1), M (6,1), P (4,2), T (3,3), S (1,3), A (4,5), D (6,5). Bir sonraki hamlemde A nesnesini merkeze yaklaştırmak için (4,4) konumuna taşımayı deneyeceğim ama (4,4) hücresinin boş olduğundan emin olmam gerekiyor. Sen de kendi tahtandaki nesneleri bu koordinatlara göre planlayabilirsin."} \\
\\ 
\\ 
\texttt{Bir sonraki komutun nedir?} \\
            }
        }
    }
    & \\ \\

    \theutterance \stepcounter{utterance}  
    & & & \multicolumn{4}{p{0.6\linewidth}}{
        \cellcolor[rgb]{0.9,0.9,0.9}{
            \makecell[{{p{\linewidth}}}]{
                \texttt{\tiny{[GM$\langle$P2]}}
                \texttt{SÖYLE: T nesnesini (3,3) konumuna taşıman harika. Benim tahtamdaki güncel durum: P (2,1), M (3,1), O (4,1), S (2,2), A (1,4), T (5,3), D (6,7). Senin A nesnesini (4,4) konumuna taşıma planın çok mantıklı. Ben de kendi tahtamdaki nesneleri senin koordinatlarına uydurmaya çalışacağım. Bir sonraki hamlemde T nesnesini merkeze yaklaştırmayı planlıyorum.} \\
            }
        }
    }
     \\ \\

    \theutterance \stepcounter{utterance}  
    & & & \multicolumn{2}{p{0.3\linewidth}}{
        \cellcolor[rgb]{0.9,0.9,0.9}{
            \makecell[{{p{\linewidth}}}]{
                \texttt{\tiny{[GM$|$GM]}}
                \texttt{SÖYLE: T nesnesini (3,3) konumuna taşıman harika. Benim tahtamdaki güncel durum: P (2,1), M (3,1), O (4,1), S (2,2), A (1,4), T (5,3), D (6,7). Senin A nesnesini (4,4) konumuna taşıma planın çok mantıklı. Ben de kendi tahtamdaki nesneleri senin koordinatlarına uydurmaya çalışacağım. Bir sonraki hamlemde T nesnesini merkeze yaklaştırmayı planlıyorum.} \\
            }
        }
    }
    & & \\ \\

    \theutterance \stepcounter{utterance}  
    & & \multicolumn{4}{p{0.6\linewidth}}{
        \cellcolor[rgb]{0.9,0.9,0.9}{
            \makecell[{{p{\linewidth}}}]{
                \texttt{\tiny{[P1$\langle$GM]}}
                \texttt{Mesajın diğer oyuncuya iletildi.} \\
\\ 
\\ 
\texttt{Şu anda en fazla 28 turdan 16. turu oynuyorsunuz.} \\
\texttt{Toplamda 16 cezanın 16 kadarını biriktirdiniz.Diğer oyuncu şu mesajı gönderdi:} \\
\\ 
\\ 
\texttt{"T nesnesini (3,3) konumuna taşıman harika. Benim tahtamdaki güncel durum: P (2,1), M (3,1), O (4,1), S (2,2), A (1,4), T (5,3), D (6,7). Senin A nesnesini (4,4) konumuna taşıma planın çok mantıklı. Ben de kendi tahtamdaki nesneleri senin koordinatlarına uydurmaya çalışacağım. Bir sonraki hamlemde T nesnesini merkeze yaklaştırmayı planlıyorum."} \\
\\ 
\\ 
\texttt{Bir sonraki komutun nedir?} \\
            }
        }
    }
    & \\ \\

    \theutterance \stepcounter{utterance}  
    & \multicolumn{4}{p{0.6\linewidth}}{
        \cellcolor[rgb]{0.9,0.9,0.9}{
            \makecell[{{p{\linewidth}}}]{
                \texttt{\tiny{[P1$\rangle$GM]}}
                \texttt{SÖYLE: A nesnesini (4,4) konumuna taşımayı planlıyorum ancak tahtayı tekrar kontrol ettiğimde (4,4) hücresinin çerçeve çizgisi (┼) içerdiğini fark ettim. Bu yüzden A nesnesini (4,4) yerine, tamamen boş olan (3,4) veya (4,3) gibi bir hücreye taşımalıyım. Bir sonraki hamlemde A nesnesini (4,3) konumuna taşımayı deneyeceğim. Sen de kendi tahtandaki T nesnesini merkeze yaklaştırırken dikkatli ol.} \\
            }
        }
    }
    & & \\ \\

    \theutterance \stepcounter{utterance}  
    & & & \multicolumn{2}{p{0.3\linewidth}}{
        \cellcolor[rgb]{0.9,0.9,0.9}{
            \makecell[{{p{\linewidth}}}]{
                \texttt{\tiny{[GM$|$GM]}}
                \texttt{SÖYLE: A nesnesini (4,4) konumuna taşımayı planlıyorum ancak tahtayı tekrar kontrol ettiğimde (4,4) hücresinin çerçeve çizgisi (┼) içerdiğini fark ettim. Bu yüzden A nesnesini (4,4) yerine, tamamen boş olan (3,4) veya (4,3) gibi bir hücreye taşımalıyım. Bir sonraki hamlemde A nesnesini (4,3) konumuna taşımayı deneyeceğim. Sen de kendi tahtandaki T nesnesini merkeze yaklaştırırken dikkatli ol.} \\
            }
        }
    }
    & & \\ \\

    \theutterance \stepcounter{utterance}  
    & & \multicolumn{4}{p{0.6\linewidth}}{
        \cellcolor[rgb]{0.9,0.9,0.9}{
            \makecell[{{p{\linewidth}}}]{
                \texttt{\tiny{[GM$\rangle$P2]}}
                \texttt{Mesajın diğer oyuncuya iletildi.} \\
\\ 
\\ 
\texttt{Şu anda en fazla 28 turdan 17. turu oynuyorsunuz.} \\
\texttt{Toplamda 16 cezanın 16 kadarını biriktirdiniz.Diğer oyuncu şu mesajı gönderdi:} \\
\\ 
\\ 
\texttt{"A nesnesini (4,4) konumuna taşımayı planlıyorum ancak tahtayı tekrar kontrol ettiğimde (4,4) hücresinin çerçeve çizgisi (┼) içerdiğini fark ettim. Bu yüzden A nesnesini (4,4) yerine, tamamen boş olan (3,4) veya (4,3) gibi bir hücreye taşımalıyım. Bir sonraki hamlemde A nesnesini (4,3) konumuna taşımayı deneyeceğim. Sen de kendi tahtandaki T nesnesini merkeze yaklaştırırken dikkatli ol."} \\
\\ 
\\ 
\texttt{Bir sonraki komutun nedir?} \\
            }
        }
    }
    & \\ \\

    \theutterance \stepcounter{utterance}  
    & & & \multicolumn{4}{p{0.6\linewidth}}{
        \cellcolor[rgb]{0.9,0.9,0.9}{
            \makecell[{{p{\linewidth}}}]{
                \texttt{\tiny{[GM$\langle$P2]}}
                \texttt{HAMLE: T, (4,3)} \\
            }
        }
    }
     \\ \\

    \theutterance \stepcounter{utterance}  
    & & & \multicolumn{2}{p{0.3\linewidth}}{
        \cellcolor[rgb]{0.9,0.9,0.9}{
            \makecell[{{p{\linewidth}}}]{
                \texttt{\tiny{[GM$|$GM]}}
                \texttt{HAMLE: T, (4,3)} \\
            }
        }
    }
    & & \\ \\

    \theutterance \stepcounter{utterance}  
    & & & \multicolumn{2}{p{0.3\linewidth}}{
        \cellcolor[rgb]{0.9,0.9,0.9}{
            \makecell[{{p{\linewidth}}}]{
                \texttt{\tiny{[GM$|$GM]}}
                \texttt{T moved to (4, 3)} \\
            }
        }
    }
    & & \\ \\

    \theutterance \stepcounter{utterance}  
    & & \multicolumn{4}{p{0.6\linewidth}}{
        \cellcolor[rgb]{0.9,0.9,0.9}{
            \makecell[{{p{\linewidth}}}]{
                \texttt{\tiny{[P1$\langle$GM]}}
                \texttt{Mesajın diğer oyuncuya iletildi.} \\
\\ 
\\ 
\texttt{Şu anda en fazla 28 turdan 17. turu oynuyorsunuz.} \\
\texttt{Toplamda 16 cezanın 16 kadarını biriktirdiniz.Diğer oyuncu kendi ızgarasında bir nesneyi hareket ettirdi.} \\
\\ 
\\ 
\texttt{Bir sonraki komutun nedir?} \\
            }
        }
    }
    & \\ \\

    \theutterance \stepcounter{utterance}  
    & \multicolumn{4}{p{0.6\linewidth}}{
        \cellcolor[rgb]{0.9,0.9,0.9}{
            \makecell[{{p{\linewidth}}}]{
                \texttt{\tiny{[P1$\rangle$GM]}}
                \texttt{SÖYLE: Diğer oyuncu T nesnesini nereye taşıdı? Ben de bir sonraki hamlemde A nesnesini (4,3) konumuna taşımayı deneyeceğim ama hata yapmamak için hücrenin tamamen boş olduğundan emin olmam gerekiyor. Lütfen onun hamlesini paylaşır mısın?} \\
            }
        }
    }
    & & \\ \\

    \theutterance \stepcounter{utterance}  
    & & & \multicolumn{2}{p{0.3\linewidth}}{
        \cellcolor[rgb]{0.9,0.9,0.9}{
            \makecell[{{p{\linewidth}}}]{
                \texttt{\tiny{[GM$|$GM]}}
                \texttt{SÖYLE: Diğer oyuncu T nesnesini nereye taşıdı? Ben de bir sonraki hamlemde A nesnesini (4,3) konumuna taşımayı deneyeceğim ama hata yapmamak için hücrenin tamamen boş olduğundan emin olmam gerekiyor. Lütfen onun hamlesini paylaşır mısın?} \\
            }
        }
    }
    & & \\ \\

    \theutterance \stepcounter{utterance}  
    & & \multicolumn{4}{p{0.6\linewidth}}{
        \cellcolor[rgb]{0.9,0.9,0.9}{
            \makecell[{{p{\linewidth}}}]{
                \texttt{\tiny{[GM$\rangle$P2]}}
                \texttt{T’ başarıyla (4,3) koordinatlarına taşındı. Güncellenmiş ızgaran şu şekilde:} \\
\texttt{```} \\
\texttt{ 1234567} \\
\texttt{╔════╤══╗} \\
\texttt{║◌PMO│◌◌║ 1} \\
\texttt{║◌S◌◌├──╢ 2} \\
\texttt{║◌◌◌T│◌◌║ 3} \\
\texttt{║A┌──┼──╢ 4} \\
\texttt{║◌│◌◌│◌◌║ 5} \\
\texttt{╟─┴┬─┤◌◌║ 6} \\
\texttt{║◌◌│◌│D◌║ 7} \\
\texttt{╚══╧═╧══╝} \\
\texttt{```} \\
\texttt{Şu anda en fazla 28 turdan 18. turu oynuyorsunuz.} \\
\texttt{Toplamda 16 cezanın 16 kadarını biriktirdiniz.Diğer oyuncu şu mesajı gönderdi:} \\
\\ 
\\ 
\texttt{"Diğer oyuncu T nesnesini nereye taşıdı? Ben de bir sonraki hamlemde A nesnesini (4,3) konumuna taşımayı deneyeceğim ama hata yapmamak için hücrenin tamamen boş olduğundan emin olmam gerekiyor. Lütfen onun hamlesini paylaşır mısın?"} \\
\\ 
\\ 
\texttt{Bir sonraki komutun nedir?} \\
            }
        }
    }
    & \\ \\

    \theutterance \stepcounter{utterance}  
    & & & \multicolumn{4}{p{0.6\linewidth}}{
        \cellcolor[rgb]{0.9,0.9,0.9}{
            \makecell[{{p{\linewidth}}}]{
                \texttt{\tiny{[GM$\langle$P2]}}
                \texttt{SÖYLE: Ben T nesnesini (4,3) konumuna taşıdım. Dikkat et, benim tahtamda (4,3) hücresi şu an T nesnesi ile dolu. Eğer senin tahtanda da (4,3) hücresi doluysa, A nesnesini oraya taşıyamazsın. Planını tekrar kontrol etmeni öneririm.} \\
            }
        }
    }
     \\ \\

    \theutterance \stepcounter{utterance}  
    & & & \multicolumn{2}{p{0.3\linewidth}}{
        \cellcolor[rgb]{0.9,0.9,0.9}{
            \makecell[{{p{\linewidth}}}]{
                \texttt{\tiny{[GM$|$GM]}}
                \texttt{SÖYLE: Ben T nesnesini (4,3) konumuna taşıdım. Dikkat et, benim tahtamda (4,3) hücresi şu an T nesnesi ile dolu. Eğer senin tahtanda da (4,3) hücresi doluysa, A nesnesini oraya taşıyamazsın. Planını tekrar kontrol etmeni öneririm.} \\
            }
        }
    }
    & & \\ \\

    \theutterance \stepcounter{utterance}  
    & & \multicolumn{4}{p{0.6\linewidth}}{
        \cellcolor[rgb]{0.9,0.9,0.9}{
            \makecell[{{p{\linewidth}}}]{
                \texttt{\tiny{[P1$\langle$GM]}}
                \texttt{Mesajın diğer oyuncuya iletildi.} \\
\\ 
\\ 
\texttt{Şu anda en fazla 28 turdan 18. turu oynuyorsunuz.} \\
\texttt{Toplamda 16 cezanın 16 kadarını biriktirdiniz.Diğer oyuncu şu mesajı gönderdi:} \\
\\ 
\\ 
\texttt{"Ben T nesnesini (4,3) konumuna taşıdım. Dikkat et, benim tahtamda (4,3) hücresi şu an T nesnesi ile dolu. Eğer senin tahtanda da (4,3) hücresi doluysa, A nesnesini oraya taşıyamazsın. Planını tekrar kontrol etmeni öneririm."} \\
\\ 
\\ 
\texttt{Bir sonraki komutun nedir?} \\
            }
        }
    }
    & \\ \\

    \theutterance \stepcounter{utterance}  
    & \multicolumn{4}{p{0.6\linewidth}}{
        \cellcolor[rgb]{0.9,0.9,0.9}{
            \makecell[{{p{\linewidth}}}]{
                \texttt{\tiny{[P1$\rangle$GM]}}
                \texttt{SÖYLE: Uyarın için teşekkürler! Kendi tahtamı tekrar kontrol ettim. (4,3) hücresi şu an boş görünüyor, ancak senin tahtandaki T nesnesiyle çakışmaması için daha güvenli bir yer seçmeliyim. A nesnesini (3,4) konumuna taşımaya karar verdim. (3,4) hücresinin tamamen boş olduğundan emin oldum.} \\
            }
        }
    }
    & & \\ \\

    \theutterance \stepcounter{utterance}  
    & & & \multicolumn{2}{p{0.3\linewidth}}{
        \cellcolor[rgb]{0.9,0.9,0.9}{
            \makecell[{{p{\linewidth}}}]{
                \texttt{\tiny{[GM$|$GM]}}
                \texttt{SÖYLE: Uyarın için teşekkürler! Kendi tahtamı tekrar kontrol ettim. (4,3) hücresi şu an boş görünüyor, ancak senin tahtandaki T nesnesiyle çakışmaması için daha güvenli bir yer seçmeliyim. A nesnesini (3,4) konumuna taşımaya karar verdim. (3,4) hücresinin tamamen boş olduğundan emin oldum.} \\
            }
        }
    }
    & & \\ \\

    \theutterance \stepcounter{utterance}  
    & & \multicolumn{4}{p{0.6\linewidth}}{
        \cellcolor[rgb]{0.9,0.9,0.9}{
            \makecell[{{p{\linewidth}}}]{
                \texttt{\tiny{[GM$\rangle$P2]}}
                \texttt{Mesajın diğer oyuncuya iletildi.} \\
\\ 
\\ 
\texttt{Şu anda en fazla 28 turdan 19. turu oynuyorsunuz.} \\
\texttt{Toplamda 16 cezanın 16 kadarını biriktirdiniz.Diğer oyuncu şu mesajı gönderdi:} \\
\\ 
\\ 
\texttt{"Uyarın için teşekkürler! Kendi tahtamı tekrar kontrol ettim. (4,3) hücresi şu an boş görünüyor, ancak senin tahtandaki T nesnesiyle çakışmaması için daha güvenli bir yer seçmeliyim. A nesnesini (3,4) konumuna taşımaya karar verdim. (3,4) hücresinin tamamen boş olduğundan emin oldum."} \\
\\ 
\\ 
\texttt{Bir sonraki komutun nedir?} \\
            }
        }
    }
    & \\ \\

    \theutterance \stepcounter{utterance}  
    & & & \multicolumn{4}{p{0.6\linewidth}}{
        \cellcolor[rgb]{0.9,0.9,0.9}{
            \makecell[{{p{\linewidth}}}]{
                \texttt{\tiny{[GM$\langle$P2]}}
                \texttt{SÖYLE: A nesnesini (3,4) konumuna taşıma kararın çok yerinde, böylece çakışmayı önlemiş olduk. Benim tahtamdaki güncel durum: P (2,1), M (3,1), O (4,1), S (2,2), T (4,3), A (1,4), D (6,7). Bir sonraki hamlemde A nesnesini merkeze yaklaştırmak için hareket ettirmeyi planlıyorum.} \\
            }
        }
    }
     \\ \\

    \theutterance \stepcounter{utterance}  
    & & & \multicolumn{2}{p{0.3\linewidth}}{
        \cellcolor[rgb]{0.9,0.9,0.9}{
            \makecell[{{p{\linewidth}}}]{
                \texttt{\tiny{[GM$|$GM]}}
                \texttt{SÖYLE: A nesnesini (3,4) konumuna taşıma kararın çok yerinde, böylece çakışmayı önlemiş olduk. Benim tahtamdaki güncel durum: P (2,1), M (3,1), O (4,1), S (2,2), T (4,3), A (1,4), D (6,7). Bir sonraki hamlemde A nesnesini merkeze yaklaştırmak için hareket ettirmeyi planlıyorum.} \\
            }
        }
    }
    & & \\ \\

    \theutterance \stepcounter{utterance}  
    & & \multicolumn{4}{p{0.6\linewidth}}{
        \cellcolor[rgb]{0.9,0.9,0.9}{
            \makecell[{{p{\linewidth}}}]{
                \texttt{\tiny{[P1$\langle$GM]}}
                \texttt{Mesajın diğer oyuncuya iletildi.} \\
\\ 
\\ 
\texttt{Şu anda en fazla 28 turdan 19. turu oynuyorsunuz.} \\
\texttt{Toplamda 16 cezanın 16 kadarını biriktirdiniz.Diğer oyuncu şu mesajı gönderdi:} \\
\\ 
\\ 
\texttt{"A nesnesini (3,4) konumuna taşıma kararın çok yerinde, böylece çakışmayı önlemiş olduk. Benim tahtamdaki güncel durum: P (2,1), M (3,1), O (4,1), S (2,2), T (4,3), A (1,4), D (6,7). Bir sonraki hamlemde A nesnesini merkeze yaklaştırmak için hareket ettirmeyi planlıyorum."} \\
\\ 
\\ 
\texttt{Bir sonraki komutun nedir?} \\
            }
        }
    }
    & \\ \\

    \theutterance \stepcounter{utterance}  
    & \multicolumn{4}{p{0.6\linewidth}}{
        \cellcolor[rgb]{0.9,0.9,0.9}{
            \makecell[{{p{\linewidth}}}]{
                \texttt{\tiny{[P1$\rangle$GM]}}
                \texttt{HAMLE: A, (3, 4)} \\
            }
        }
    }
    & & \\ \\

    \theutterance \stepcounter{utterance}  
    & & & \multicolumn{2}{p{0.3\linewidth}}{
        \cellcolor[rgb]{0.9,0.9,0.9}{
            \makecell[{{p{\linewidth}}}]{
                \texttt{\tiny{[GM$|$GM]}}
                \texttt{HAMLE: A, (3, 4)} \\
            }
        }
    }
    & & \\ \\

    \theutterance \stepcounter{utterance}  
    & & & \multicolumn{2}{p{0.3\linewidth}}{
        \cellcolor[rgb]{0.9,0.9,0.9}{
            \makecell[{{p{\linewidth}}}]{
                \texttt{\tiny{[GM$|$GM]}}
                \texttt{Ceza: (3,4) boş değil; "─" içeriyor.} \\
\\ 
\texttt{Toplamda 16 cezanın 17 kadarını biriktirdiniz.Lütfen tekrar deneyin!} \\
            }
        }
    }
    & & \\ \\

    \theutterance \stepcounter{utterance}  
    & & \multicolumn{4}{p{0.6\linewidth}}{
        \cellcolor[rgb]{0.9,0.9,0.9}{
            \makecell[{{p{\linewidth}}}]{
                \texttt{\tiny{[P1$\langle$GM]}}
                \texttt{Ceza: (3,4) boş değil; "─" içeriyor.} \\
\\ 
\texttt{Toplamda 16 cezanın 17 kadarını biriktirdiniz.Lütfen tekrar deneyin!} \\
            }
        }
    }
    & \\ \\

    \theutterance \stepcounter{utterance}  
    & \multicolumn{4}{p{0.6\linewidth}}{
        \cellcolor[rgb]{0.9,0.9,0.9}{
            \makecell[{{p{\linewidth}}}]{
                \texttt{\tiny{[P1$\rangle$GM]}}
                \texttt{HAMLE: A, (3, 5)} \\
            }
        }
    }
    & & \\ \\

    \theutterance \stepcounter{utterance}  
    & & & \multicolumn{2}{p{0.3\linewidth}}{
        \cellcolor[rgb]{0.9,0.9,0.9}{
            \makecell[{{p{\linewidth}}}]{
                \texttt{\tiny{[GM$|$GM]}}
                \texttt{HAMLE: A, (3, 5)} \\
            }
        }
    }
    & & \\ \\

    \theutterance \stepcounter{utterance}  
    & & & \multicolumn{2}{p{0.3\linewidth}}{
        \cellcolor[rgb]{0.9,0.9,0.9}{
            \makecell[{{p{\linewidth}}}]{
                \texttt{\tiny{[GM$|$GM]}}
                \texttt{A moved to (3, 5)} \\
            }
        }
    }
    & & \\ \\

    \theutterance \stepcounter{utterance}  
    & & \multicolumn{4}{p{0.6\linewidth}}{
        \cellcolor[rgb]{0.9,0.9,0.9}{
            \makecell[{{p{\linewidth}}}]{
                \texttt{\tiny{[GM$\rangle$P2]}}
                \texttt{Mesajın diğer oyuncuya iletildi.} \\
\\ 
\\ 
\texttt{Şu anda en fazla 28 turdan 20. turu oynuyorsunuz.} \\
\texttt{Toplamda 16 cezanın 17 kadarını biriktirdiniz.Diğer oyuncu kendi ızgarasında bir nesneyi hareket ettirdi.} \\
\\ 
\\ 
\texttt{Bir sonraki komutun nedir?} \\
            }
        }
    }
    & \\ \\

    \theutterance \stepcounter{utterance}  
    & & & \multicolumn{4}{p{0.6\linewidth}}{
        \cellcolor[rgb]{0.9,0.9,0.9}{
            \makecell[{{p{\linewidth}}}]{
                \texttt{\tiny{[GM$\langle$P2]}}
                \texttt{SÖYLE: Ben A nesnesini merkeze yaklaştırmayı planlıyordum. Sen A nesnesini (3,4) konumuna başarıyla taşıdın mı? Kendi tahtandaki güncel durumu ve bir sonraki hamleni paylaşırsan ben de ona göre hareket edebilirim.} \\
            }
        }
    }
     \\ \\

    \theutterance \stepcounter{utterance}  
    & & & \multicolumn{2}{p{0.3\linewidth}}{
        \cellcolor[rgb]{0.9,0.9,0.9}{
            \makecell[{{p{\linewidth}}}]{
                \texttt{\tiny{[GM$|$GM]}}
                \texttt{SÖYLE: Ben A nesnesini merkeze yaklaştırmayı planlıyordum. Sen A nesnesini (3,4) konumuna başarıyla taşıdın mı? Kendi tahtandaki güncel durumu ve bir sonraki hamleni paylaşırsan ben de ona göre hareket edebilirim.} \\
            }
        }
    }
    & & \\ \\

    \theutterance \stepcounter{utterance}  
    & & \multicolumn{4}{p{0.6\linewidth}}{
        \cellcolor[rgb]{0.9,0.9,0.9}{
            \makecell[{{p{\linewidth}}}]{
                \texttt{\tiny{[P1$\langle$GM]}}
                \texttt{A’ başarıyla (3,5) koordinatlarına taşındı. Güncellenmiş ızgaran şu şekilde:} \\
\texttt{```} \\
\texttt{ 1234567} \\
\texttt{╔════╤══╗} \\
\texttt{║◌◌◌O│M◌║ 1} \\
\texttt{║◌◌◌P├──╢ 2} \\
\texttt{║S◌T◌│◌◌║ 3} \\
\texttt{║◌┌──┼──╢ 4} \\
\texttt{║◌│A◌│D◌║ 5} \\
\texttt{╟─┴┬─┤◌◌║ 6} \\
\texttt{║◌◌│◌│◌◌║ 7} \\
\texttt{╚══╧═╧══╝} \\
\texttt{```} \\
\texttt{Şu anda en fazla 28 turdan 20. turu oynuyorsunuz.} \\
\texttt{Toplamda 16 cezanın 17 kadarını biriktirdiniz.Diğer oyuncu şu mesajı gönderdi:} \\
\\ 
\\ 
\texttt{"Ben A nesnesini merkeze yaklaştırmayı planlıyordum. Sen A nesnesini (3,4) konumuna başarıyla taşıdın mı? Kendi tahtandaki güncel durumu ve bir sonraki hamleni paylaşırsan ben de ona göre hareket edebilirim."} \\
\\ 
\\ 
\texttt{Bir sonraki komutun nedir?} \\
            }
        }
    }
    & \\ \\

    \theutterance \stepcounter{utterance}  
    & \multicolumn{4}{p{0.6\linewidth}}{
        \cellcolor[rgb]{0.9,0.9,0.9}{
            \makecell[{{p{\linewidth}}}]{
                \texttt{\tiny{[P1$\rangle$GM]}}
                \texttt{SÖYLE: A nesnesini (3,5) konumuna başarıyla taşıdım. Benim tahtamdaki güncel durum: O (4,1), M (6,1), P (4,2), T (3,3), S (1,3), A (3,5), D (6,5). Bir sonraki hamlemde S nesnesini (2,3) konumuna çekerek merkeze daha da yaklaştırmayı planlıyorum. Sen de kendi tahtandaki A nesnesini merkeze yaklaştırabilirsin.} \\
            }
        }
    }
    & & \\ \\

    \theutterance \stepcounter{utterance}  
    & & & \multicolumn{2}{p{0.3\linewidth}}{
        \cellcolor[rgb]{0.9,0.9,0.9}{
            \makecell[{{p{\linewidth}}}]{
                \texttt{\tiny{[GM$|$GM]}}
                \texttt{SÖYLE: A nesnesini (3,5) konumuna başarıyla taşıdım. Benim tahtamdaki güncel durum: O (4,1), M (6,1), P (4,2), T (3,3), S (1,3), A (3,5), D (6,5). Bir sonraki hamlemde S nesnesini (2,3) konumuna çekerek merkeze daha da yaklaştırmayı planlıyorum. Sen de kendi tahtandaki A nesnesini merkeze yaklaştırabilirsin.} \\
            }
        }
    }
    & & \\ \\

    \theutterance \stepcounter{utterance}  
    & & \multicolumn{4}{p{0.6\linewidth}}{
        \cellcolor[rgb]{0.9,0.9,0.9}{
            \makecell[{{p{\linewidth}}}]{
                \texttt{\tiny{[GM$\rangle$P2]}}
                \texttt{Mesajın diğer oyuncuya iletildi.} \\
\\ 
\\ 
\texttt{Şu anda en fazla 28 turdan 21. turu oynuyorsunuz.} \\
\texttt{Toplamda 16 cezanın 17 kadarını biriktirdiniz.Diğer oyuncu şu mesajı gönderdi:} \\
\\ 
\\ 
\texttt{"A nesnesini (3,5) konumuna başarıyla taşıdım. Benim tahtamdaki güncel durum: O (4,1), M (6,1), P (4,2), T (3,3), S (1,3), A (3,5), D (6,5). Bir sonraki hamlemde S nesnesini (2,3) konumuna çekerek merkeze daha da yaklaştırmayı planlıyorum. Sen de kendi tahtandaki A nesnesini merkeze yaklaştırabilirsin."} \\
\\ 
\\ 
\texttt{Bir sonraki komutun nedir?} \\
            }
        }
    }
    & \\ \\

    \theutterance \stepcounter{utterance}  
    & & & \multicolumn{4}{p{0.6\linewidth}}{
        \cellcolor[rgb]{0.9,0.9,0.9}{
            \makecell[{{p{\linewidth}}}]{
                \texttt{\tiny{[GM$\langle$P2]}}
                \texttt{HAMLE: A, (2,4)} \\
            }
        }
    }
     \\ \\

    \theutterance \stepcounter{utterance}  
    & & & \multicolumn{2}{p{0.3\linewidth}}{
        \cellcolor[rgb]{0.9,0.9,0.9}{
            \makecell[{{p{\linewidth}}}]{
                \texttt{\tiny{[GM$|$GM]}}
                \texttt{HAMLE: A, (2,4)} \\
            }
        }
    }
    & & \\ \\

    \theutterance \stepcounter{utterance}  
    & & & \multicolumn{2}{p{0.3\linewidth}}{
        \cellcolor[rgb]{0.9,0.9,0.9}{
            \makecell[{{p{\linewidth}}}]{
                \texttt{\tiny{[GM$|$GM]}}
                \texttt{Ceza: (2,4) boş değil; "┌" içeriyor.} \\
\\ 
\texttt{Toplamda 16 cezanın 18 kadarını biriktirdiniz.Lütfen tekrar deneyin!} \\
            }
        }
    }
    & & \\ \\

    \theutterance \stepcounter{utterance}  
    & & \multicolumn{4}{p{0.6\linewidth}}{
        \cellcolor[rgb]{0.9,0.9,0.9}{
            \makecell[{{p{\linewidth}}}]{
                \texttt{\tiny{[GM$\rangle$P2]}}
                \texttt{Ceza: (2,4) boş değil; "┌" içeriyor.} \\
\\ 
\texttt{Toplamda 16 cezanın 18 kadarını biriktirdiniz.Lütfen tekrar deneyin!} \\
            }
        }
    }
    & \\ \\

    \theutterance \stepcounter{utterance}  
    & & & \multicolumn{4}{p{0.6\linewidth}}{
        \cellcolor[rgb]{0.9,0.9,0.9}{
            \makecell[{{p{\linewidth}}}]{
                \texttt{\tiny{[GM$\langle$P2]}}
                \texttt{HAMLE: A, (1,3)} \\
            }
        }
    }
     \\ \\

    \theutterance \stepcounter{utterance}  
    & & & \multicolumn{2}{p{0.3\linewidth}}{
        \cellcolor[rgb]{0.9,0.9,0.9}{
            \makecell[{{p{\linewidth}}}]{
                \texttt{\tiny{[GM$|$GM]}}
                \texttt{HAMLE: A, (1,3)} \\
            }
        }
    }
    & & \\ \\

    \theutterance \stepcounter{utterance}  
    & & & \multicolumn{2}{p{0.3\linewidth}}{
        \cellcolor[rgb]{0.9,0.9,0.9}{
            \makecell[{{p{\linewidth}}}]{
                \texttt{\tiny{[GM$|$GM]}}
                \texttt{A moved to (1, 3)} \\
            }
        }
    }
    & & \\ \\

    \theutterance \stepcounter{utterance}  
    & & \multicolumn{4}{p{0.6\linewidth}}{
        \cellcolor[rgb]{0.9,0.9,0.9}{
            \makecell[{{p{\linewidth}}}]{
                \texttt{\tiny{[P1$\langle$GM]}}
                \texttt{Mesajın diğer oyuncuya iletildi.} \\
\\ 
\\ 
\texttt{Şu anda en fazla 28 turdan 21. turu oynuyorsunuz.} \\
\texttt{Toplamda 16 cezanın 18 kadarını biriktirdiniz.Diğer oyuncu kendi ızgarasında bir nesneyi hareket ettirdi.} \\
\\ 
\\ 
\texttt{Bir sonraki komutun nedir?} \\
            }
        }
    }
    & \\ \\

    \theutterance \stepcounter{utterance}  
    & \multicolumn{4}{p{0.6\linewidth}}{
        \cellcolor[rgb]{0.9,0.9,0.9}{
            \makecell[{{p{\linewidth}}}]{
                \texttt{\tiny{[P1$\rangle$GM]}}
                \texttt{SÖYLE: Diğer oyuncu A nesnesini nereye taşıdı? Ben de bir sonraki hamlemde S nesnesini (2,3) konumuna taşımayı planlıyorum. (2,3) hücresinin tamamen boş olduğundan emin olmam gerekiyor. Lütfen onun hamlesini paylaşır mısın?} \\
            }
        }
    }
    & & \\ \\

    \theutterance \stepcounter{utterance}  
    & & & \multicolumn{2}{p{0.3\linewidth}}{
        \cellcolor[rgb]{0.9,0.9,0.9}{
            \makecell[{{p{\linewidth}}}]{
                \texttt{\tiny{[GM$|$GM]}}
                \texttt{SÖYLE: Diğer oyuncu A nesnesini nereye taşıdı? Ben de bir sonraki hamlemde S nesnesini (2,3) konumuna taşımayı planlıyorum. (2,3) hücresinin tamamen boş olduğundan emin olmam gerekiyor. Lütfen onun hamlesini paylaşır mısın?} \\
            }
        }
    }
    & & \\ \\

    \theutterance \stepcounter{utterance}  
    & & \multicolumn{4}{p{0.6\linewidth}}{
        \cellcolor[rgb]{0.9,0.9,0.9}{
            \makecell[{{p{\linewidth}}}]{
                \texttt{\tiny{[GM$\rangle$P2]}}
                \texttt{A’ başarıyla (1,3) koordinatlarına taşındı. Güncellenmiş ızgaran şu şekilde:} \\
\texttt{```} \\
\texttt{ 1234567} \\
\texttt{╔════╤══╗} \\
\texttt{║◌PMO│◌◌║ 1} \\
\texttt{║◌S◌◌├──╢ 2} \\
\texttt{║A◌◌T│◌◌║ 3} \\
\texttt{║◌┌──┼──╢ 4} \\
\texttt{║◌│◌◌│◌◌║ 5} \\
\texttt{╟─┴┬─┤◌◌║ 6} \\
\texttt{║◌◌│◌│D◌║ 7} \\
\texttt{╚══╧═╧══╝} \\
\texttt{```} \\
\texttt{Şu anda en fazla 28 turdan 22. turu oynuyorsunuz.} \\
\texttt{Toplamda 16 cezanın 18 kadarını biriktirdiniz.Diğer oyuncu şu mesajı gönderdi:} \\
\\ 
\\ 
\texttt{"Diğer oyuncu A nesnesini nereye taşıdı? Ben de bir sonraki hamlemde S nesnesini (2,3) konumuna taşımayı planlıyorum. (2,3) hücresinin tamamen boş olduğundan emin olmam gerekiyor. Lütfen onun hamlesini paylaşır mısın?"} \\
\\ 
\\ 
\texttt{Bir sonraki komutun nedir?} \\
            }
        }
    }
    & \\ \\

    \theutterance \stepcounter{utterance}  
    & & & \multicolumn{4}{p{0.6\linewidth}}{
        \cellcolor[rgb]{0.9,0.9,0.9}{
            \makecell[{{p{\linewidth}}}]{
                \texttt{\tiny{[GM$\langle$P2]}}
                \texttt{SÖYLE: Ben A nesnesini (1,3) konumuna taşıdım. Benim tahtamdaki güncel durum: P (2,1), M (3,1), O (4,1), S (2,2), T (4,3), A (1,3), D (6,7). S nesnesini (2,3) konumuna taşıma planın gayet iyi, umarım senin tahtanda o hücre tamamen boştur. Ben de bir sonraki hamlemde D nesnesini merkeze yaklaştırmayı planlıyorum.} \\
            }
        }
    }
     \\ \\

    \theutterance \stepcounter{utterance}  
    & & & \multicolumn{2}{p{0.3\linewidth}}{
        \cellcolor[rgb]{0.9,0.9,0.9}{
            \makecell[{{p{\linewidth}}}]{
                \texttt{\tiny{[GM$|$GM]}}
                \texttt{SÖYLE: Ben A nesnesini (1,3) konumuna taşıdım. Benim tahtamdaki güncel durum: P (2,1), M (3,1), O (4,1), S (2,2), T (4,3), A (1,3), D (6,7). S nesnesini (2,3) konumuna taşıma planın gayet iyi, umarım senin tahtanda o hücre tamamen boştur. Ben de bir sonraki hamlemde D nesnesini merkeze yaklaştırmayı planlıyorum.} \\
            }
        }
    }
    & & \\ \\

    \theutterance \stepcounter{utterance}  
    & & \multicolumn{4}{p{0.6\linewidth}}{
        \cellcolor[rgb]{0.9,0.9,0.9}{
            \makecell[{{p{\linewidth}}}]{
                \texttt{\tiny{[P1$\langle$GM]}}
                \texttt{Mesajın diğer oyuncuya iletildi.} \\
\\ 
\\ 
\texttt{Şu anda en fazla 28 turdan 22. turu oynuyorsunuz.} \\
\texttt{Toplamda 16 cezanın 18 kadarını biriktirdiniz.Diğer oyuncu şu mesajı gönderdi:} \\
\\ 
\\ 
\texttt{"Ben A nesnesini (1,3) konumuna taşıdım. Benim tahtamdaki güncel durum: P (2,1), M (3,1), O (4,1), S (2,2), T (4,3), A (1,3), D (6,7). S nesnesini (2,3) konumuna taşıma planın gayet iyi, umarım senin tahtanda o hücre tamamen boştur. Ben de bir sonraki hamlemde D nesnesini merkeze yaklaştırmayı planlıyorum."} \\
\\ 
\\ 
\texttt{Bir sonraki komutun nedir?} \\
            }
        }
    }
    & \\ \\

    \theutterance \stepcounter{utterance}  
    & \multicolumn{4}{p{0.6\linewidth}}{
        \cellcolor[rgb]{0.9,0.9,0.9}{
            \makecell[{{p{\linewidth}}}]{
                \texttt{\tiny{[P1$\rangle$GM]}}
                \texttt{HAMLE: S, (2, 3)} \\
            }
        }
    }
    & & \\ \\

    \theutterance \stepcounter{utterance}  
    & & & \multicolumn{2}{p{0.3\linewidth}}{
        \cellcolor[rgb]{0.9,0.9,0.9}{
            \makecell[{{p{\linewidth}}}]{
                \texttt{\tiny{[GM$|$GM]}}
                \texttt{HAMLE: S, (2, 3)} \\
            }
        }
    }
    & & \\ \\

    \theutterance \stepcounter{utterance}  
    & & & \multicolumn{2}{p{0.3\linewidth}}{
        \cellcolor[rgb]{0.9,0.9,0.9}{
            \makecell[{{p{\linewidth}}}]{
                \texttt{\tiny{[GM$|$GM]}}
                \texttt{S moved to (2, 3)} \\
            }
        }
    }
    & & \\ \\

    \theutterance \stepcounter{utterance}  
    & & \multicolumn{4}{p{0.6\linewidth}}{
        \cellcolor[rgb]{0.9,0.9,0.9}{
            \makecell[{{p{\linewidth}}}]{
                \texttt{\tiny{[GM$\rangle$P2]}}
                \texttt{Mesajın diğer oyuncuya iletildi.} \\
\\ 
\\ 
\texttt{Şu anda en fazla 28 turdan 23. turu oynuyorsunuz.} \\
\texttt{Toplamda 16 cezanın 18 kadarını biriktirdiniz.Diğer oyuncu kendi ızgarasında bir nesneyi hareket ettirdi.} \\
\\ 
\\ 
\texttt{Bir sonraki komutun nedir?} \\
            }
        }
    }
    & \\ \\

    \theutterance \stepcounter{utterance}  
    & & & \multicolumn{4}{p{0.6\linewidth}}{
        \cellcolor[rgb]{0.9,0.9,0.9}{
            \makecell[{{p{\linewidth}}}]{
                \texttt{\tiny{[GM$\langle$P2]}}
                \texttt{SÖYLE: Ben A nesnesini (1,3) konumuna taşıdım. Sen S nesnesini (2,3) konumuna başarıyla taşıyabildin mi? Kendi tahtandaki güncel durumu ve bir sonraki hamleni paylaşırsan ben de D nesnesini hareket ettirirken ona göre plan yapabilirim.} \\
            }
        }
    }
     \\ \\

    \theutterance \stepcounter{utterance}  
    & & & \multicolumn{2}{p{0.3\linewidth}}{
        \cellcolor[rgb]{0.9,0.9,0.9}{
            \makecell[{{p{\linewidth}}}]{
                \texttt{\tiny{[GM$|$GM]}}
                \texttt{SÖYLE: Ben A nesnesini (1,3) konumuna taşıdım. Sen S nesnesini (2,3) konumuna başarıyla taşıyabildin mi? Kendi tahtandaki güncel durumu ve bir sonraki hamleni paylaşırsan ben de D nesnesini hareket ettirirken ona göre plan yapabilirim.} \\
            }
        }
    }
    & & \\ \\

    \theutterance \stepcounter{utterance}  
    & & \multicolumn{4}{p{0.6\linewidth}}{
        \cellcolor[rgb]{0.9,0.9,0.9}{
            \makecell[{{p{\linewidth}}}]{
                \texttt{\tiny{[P1$\langle$GM]}}
                \texttt{S’ başarıyla (2,3) koordinatlarına taşındı. Güncellenmiş ızgaran şu şekilde:} \\
\texttt{```} \\
\texttt{ 1234567} \\
\texttt{╔════╤══╗} \\
\texttt{║◌◌◌O│M◌║ 1} \\
\texttt{║◌◌◌P├──╢ 2} \\
\texttt{║◌ST◌│◌◌║ 3} \\
\texttt{║◌┌──┼──╢ 4} \\
\texttt{║◌│A◌│D◌║ 5} \\
\texttt{╟─┴┬─┤◌◌║ 6} \\
\texttt{║◌◌│◌│◌◌║ 7} \\
\texttt{╚══╧═╧══╝} \\
\texttt{```} \\
\texttt{Şu anda en fazla 28 turdan 23. turu oynuyorsunuz.} \\
\texttt{Toplamda 16 cezanın 18 kadarını biriktirdiniz.Diğer oyuncu şu mesajı gönderdi:} \\
\\ 
\\ 
\texttt{"Ben A nesnesini (1,3) konumuna taşıdım. Sen S nesnesini (2,3) konumuna başarıyla taşıyabildin mi? Kendi tahtandaki güncel durumu ve bir sonraki hamleni paylaşırsan ben de D nesnesini hareket ettirirken ona göre plan yapabilirim."} \\
\\ 
\\ 
\texttt{Bir sonraki komutun nedir?} \\
            }
        }
    }
    & \\ \\

    \theutterance \stepcounter{utterance}  
    & \multicolumn{4}{p{0.6\linewidth}}{
        \cellcolor[rgb]{0.9,0.9,0.9}{
            \makecell[{{p{\linewidth}}}]{
                \texttt{\tiny{[P1$\rangle$GM]}}
                \texttt{SÖYLE: S nesnesini (2,3) konumuna başarıyla taşıdım. Benim tahtamdaki güncel durum: O (4,1), M (6,1), P (4,2), T (3,3), S (2,3), A (3,5), D (6,5). Bir sonraki hamlemde T nesnesini (3,2) konumuna çekerek daha da toplu bir hale getirmeyi planlıyorum. Sen D nesnesini hareket ettirirken dikkatli ol, hata yapmamaya çalışalım.} \\
            }
        }
    }
    & & \\ \\

    \theutterance \stepcounter{utterance}  
    & & & \multicolumn{2}{p{0.3\linewidth}}{
        \cellcolor[rgb]{0.9,0.9,0.9}{
            \makecell[{{p{\linewidth}}}]{
                \texttt{\tiny{[GM$|$GM]}}
                \texttt{SÖYLE: S nesnesini (2,3) konumuna başarıyla taşıdım. Benim tahtamdaki güncel durum: O (4,1), M (6,1), P (4,2), T (3,3), S (2,3), A (3,5), D (6,5). Bir sonraki hamlemde T nesnesini (3,2) konumuna çekerek daha da toplu bir hale getirmeyi planlıyorum. Sen D nesnesini hareket ettirirken dikkatli ol, hata yapmamaya çalışalım.} \\
            }
        }
    }
    & & \\ \\

    \theutterance \stepcounter{utterance}  
    & & \multicolumn{4}{p{0.6\linewidth}}{
        \cellcolor[rgb]{0.9,0.9,0.9}{
            \makecell[{{p{\linewidth}}}]{
                \texttt{\tiny{[GM$\rangle$P2]}}
                \texttt{Mesajın diğer oyuncuya iletildi.} \\
\\ 
\\ 
\texttt{Şu anda en fazla 28 turdan 24. turu oynuyorsunuz.} \\
\texttt{Toplamda 16 cezanın 18 kadarını biriktirdiniz.Diğer oyuncu şu mesajı gönderdi:} \\
\\ 
\\ 
\texttt{"S nesnesini (2,3) konumuna başarıyla taşıdım. Benim tahtamdaki güncel durum: O (4,1), M (6,1), P (4,2), T (3,3), S (2,3), A (3,5), D (6,5). Bir sonraki hamlemde T nesnesini (3,2) konumuna çekerek daha da toplu bir hale getirmeyi planlıyorum. Sen D nesnesini hareket ettirirken dikkatli ol, hata yapmamaya çalışalım."} \\
\\ 
\\ 
\texttt{Bir sonraki komutun nedir?} \\
            }
        }
    }
    & \\ \\

    \theutterance \stepcounter{utterance}  
    & & & \multicolumn{4}{p{0.6\linewidth}}{
        \cellcolor[rgb]{0.9,0.9,0.9}{
            \makecell[{{p{\linewidth}}}]{
                \texttt{\tiny{[GM$\langle$P2]}}
                \texttt{HAMLE: D, (6,6)} \\
            }
        }
    }
     \\ \\

    \theutterance \stepcounter{utterance}  
    & & & \multicolumn{2}{p{0.3\linewidth}}{
        \cellcolor[rgb]{0.9,0.9,0.9}{
            \makecell[{{p{\linewidth}}}]{
                \texttt{\tiny{[GM$|$GM]}}
                \texttt{HAMLE: D, (6,6)} \\
            }
        }
    }
    & & \\ \\

    \theutterance \stepcounter{utterance}  
    & & & \multicolumn{2}{p{0.3\linewidth}}{
        \cellcolor[rgb]{0.9,0.9,0.9}{
            \makecell[{{p{\linewidth}}}]{
                \texttt{\tiny{[GM$|$GM]}}
                \texttt{D moved to (6, 6)} \\
            }
        }
    }
    & & \\ \\

    \theutterance \stepcounter{utterance}  
    & & \multicolumn{4}{p{0.6\linewidth}}{
        \cellcolor[rgb]{0.9,0.9,0.9}{
            \makecell[{{p{\linewidth}}}]{
                \texttt{\tiny{[P1$\langle$GM]}}
                \texttt{Mesajın diğer oyuncuya iletildi.} \\
\\ 
\\ 
\texttt{Şu anda en fazla 28 turdan 24. turu oynuyorsunuz.} \\
\texttt{Toplamda 16 cezanın 18 kadarını biriktirdiniz.Diğer oyuncu kendi ızgarasında bir nesneyi hareket ettirdi.} \\
\\ 
\\ 
\texttt{Bir sonraki komutun nedir?} \\
            }
        }
    }
    & \\ \\

    \theutterance \stepcounter{utterance}  
    & \multicolumn{4}{p{0.6\linewidth}}{
        \cellcolor[rgb]{0.9,0.9,0.9}{
            \makecell[{{p{\linewidth}}}]{
                \texttt{\tiny{[P1$\rangle$GM]}}
                \texttt{SÖYLE: Diğer oyuncu D nesnesini nereye taşıdı? Ben de bir sonraki hamlemde T nesnesini (3,2) konumuna taşımayı planlıyorum. (3,2) hücresinin tamamen boş olduğundan emin olmam gerekiyor. Lütfen onun hamlesini paylaşır mısın?} \\
            }
        }
    }
    & & \\ \\

    \theutterance \stepcounter{utterance}  
    & & & \multicolumn{2}{p{0.3\linewidth}}{
        \cellcolor[rgb]{0.9,0.9,0.9}{
            \makecell[{{p{\linewidth}}}]{
                \texttt{\tiny{[GM$|$GM]}}
                \texttt{SÖYLE: Diğer oyuncu D nesnesini nereye taşıdı? Ben de bir sonraki hamlemde T nesnesini (3,2) konumuna taşımayı planlıyorum. (3,2) hücresinin tamamen boş olduğundan emin olmam gerekiyor. Lütfen onun hamlesini paylaşır mısın?} \\
            }
        }
    }
    & & \\ \\

    \theutterance \stepcounter{utterance}  
    & & \multicolumn{4}{p{0.6\linewidth}}{
        \cellcolor[rgb]{0.9,0.9,0.9}{
            \makecell[{{p{\linewidth}}}]{
                \texttt{\tiny{[GM$\rangle$P2]}}
                \texttt{D’ başarıyla (6,6) koordinatlarına taşındı. Güncellenmiş ızgaran şu şekilde:} \\
\texttt{```} \\
\texttt{ 1234567} \\
\texttt{╔════╤══╗} \\
\texttt{║◌PMO│◌◌║ 1} \\
\texttt{║◌S◌◌├──╢ 2} \\
\texttt{║A◌◌T│◌◌║ 3} \\
\texttt{║◌┌──┼──╢ 4} \\
\texttt{║◌│◌◌│◌◌║ 5} \\
\texttt{╟─┴┬─┤D◌║ 6} \\
\texttt{║◌◌│◌│◌◌║ 7} \\
\texttt{╚══╧═╧══╝} \\
\texttt{```} \\
\texttt{Şu anda en fazla 28 turdan 25. turu oynuyorsunuz.} \\
\texttt{Toplamda 16 cezanın 18 kadarını biriktirdiniz.Diğer oyuncu şu mesajı gönderdi:} \\
\\ 
\\ 
\texttt{"Diğer oyuncu D nesnesini nereye taşıdı? Ben de bir sonraki hamlemde T nesnesini (3,2) konumuna taşımayı planlıyorum. (3,2) hücresinin tamamen boş olduğundan emin olmam gerekiyor. Lütfen onun hamlesini paylaşır mısın?"} \\
\\ 
\\ 
\texttt{Bir sonraki komutun nedir?} \\
            }
        }
    }
    & \\ \\

    \theutterance \stepcounter{utterance}  
    & & & \multicolumn{4}{p{0.6\linewidth}}{
        \cellcolor[rgb]{0.9,0.9,0.9}{
            \makecell[{{p{\linewidth}}}]{
                \texttt{\tiny{[GM$\langle$P2]}}
                \texttt{SÖYLE: Ben D nesnesini (6,6) konumuna taşıdım. Benim tahtamdaki güncel durum: P (2,1), M (3,1), O (4,1), S (2,2), T (4,3), A (1,3), D (6,6). Senin T nesnesini (3,2) konumuna taşıma planın çok mantıklı, umarım o hücre senin tahtanda boştur. Ben de bir sonraki hamlemde D nesnesini biraz daha yukarı çekmeyi planlıyorum.} \\
            }
        }
    }
     \\ \\

    \theutterance \stepcounter{utterance}  
    & & & \multicolumn{2}{p{0.3\linewidth}}{
        \cellcolor[rgb]{0.9,0.9,0.9}{
            \makecell[{{p{\linewidth}}}]{
                \texttt{\tiny{[GM$|$GM]}}
                \texttt{SÖYLE: Ben D nesnesini (6,6) konumuna taşıdım. Benim tahtamdaki güncel durum: P (2,1), M (3,1), O (4,1), S (2,2), T (4,3), A (1,3), D (6,6). Senin T nesnesini (3,2) konumuna taşıma planın çok mantıklı, umarım o hücre senin tahtanda boştur. Ben de bir sonraki hamlemde D nesnesini biraz daha yukarı çekmeyi planlıyorum.} \\
            }
        }
    }
    & & \\ \\

    \theutterance \stepcounter{utterance}  
    & & \multicolumn{4}{p{0.6\linewidth}}{
        \cellcolor[rgb]{0.9,0.9,0.9}{
            \makecell[{{p{\linewidth}}}]{
                \texttt{\tiny{[P1$\langle$GM]}}
                \texttt{Mesajın diğer oyuncuya iletildi.} \\
\\ 
\\ 
\texttt{Şu anda en fazla 28 turdan 25. turu oynuyorsunuz.} \\
\texttt{Toplamda 16 cezanın 18 kadarını biriktirdiniz.Diğer oyuncu şu mesajı gönderdi:} \\
\\ 
\\ 
\texttt{"Ben D nesnesini (6,6) konumuna taşıdım. Benim tahtamdaki güncel durum: P (2,1), M (3,1), O (4,1), S (2,2), T (4,3), A (1,3), D (6,6). Senin T nesnesini (3,2) konumuna taşıma planın çok mantıklı, umarım o hücre senin tahtanda boştur. Ben de bir sonraki hamlemde D nesnesini biraz daha yukarı çekmeyi planlıyorum."} \\
\\ 
\\ 
\texttt{Bir sonraki komutun nedir?} \\
            }
        }
    }
    & \\ \\

    \theutterance \stepcounter{utterance}  
    & \multicolumn{4}{p{0.6\linewidth}}{
        \cellcolor[rgb]{0.9,0.9,0.9}{
            \makecell[{{p{\linewidth}}}]{
                \texttt{\tiny{[P1$\rangle$GM]}}
                \texttt{HAMLE: T, (3, 2)} \\
            }
        }
    }
    & & \\ \\

    \theutterance \stepcounter{utterance}  
    & & & \multicolumn{2}{p{0.3\linewidth}}{
        \cellcolor[rgb]{0.9,0.9,0.9}{
            \makecell[{{p{\linewidth}}}]{
                \texttt{\tiny{[GM$|$GM]}}
                \texttt{HAMLE: T, (3, 2)} \\
            }
        }
    }
    & & \\ \\

    \theutterance \stepcounter{utterance}  
    & & & \multicolumn{2}{p{0.3\linewidth}}{
        \cellcolor[rgb]{0.9,0.9,0.9}{
            \makecell[{{p{\linewidth}}}]{
                \texttt{\tiny{[GM$|$GM]}}
                \texttt{T moved to (3, 2)} \\
            }
        }
    }
    & & \\ \\

    \theutterance \stepcounter{utterance}  
    & & \multicolumn{4}{p{0.6\linewidth}}{
        \cellcolor[rgb]{0.9,0.9,0.9}{
            \makecell[{{p{\linewidth}}}]{
                \texttt{\tiny{[GM$\rangle$P2]}}
                \texttt{Mesajın diğer oyuncuya iletildi.} \\
\\ 
\\ 
\texttt{Şu anda en fazla 28 turdan 26. turu oynuyorsunuz.} \\
\texttt{Toplamda 16 cezanın 18 kadarını biriktirdiniz.Diğer oyuncu kendi ızgarasında bir nesneyi hareket ettirdi.} \\
\\ 
\\ 
\texttt{Bir sonraki komutun nedir?} \\
            }
        }
    }
    & \\ \\

    \theutterance \stepcounter{utterance}  
    & & & \multicolumn{4}{p{0.6\linewidth}}{
        \cellcolor[rgb]{0.9,0.9,0.9}{
            \makecell[{{p{\linewidth}}}]{
                \texttt{\tiny{[GM$\langle$P2]}}
                \texttt{SÖYLE: Ben D nesnesini (6,6) konumuna taşıdım. Sen T nesnesini (3,2) konumuna başarıyla taşıyabildin mi? Kendi tahtandaki güncel durumu ve bir sonraki hamleni paylaşırsan, ben de D nesnesini (6,5) veya (6,4) gibi daha yukarıdaki bir konuma çekerek seninle uyumlu hale getirmeye çalışabilirim.} \\
            }
        }
    }
     \\ \\

    \theutterance \stepcounter{utterance}  
    & & & \multicolumn{2}{p{0.3\linewidth}}{
        \cellcolor[rgb]{0.9,0.9,0.9}{
            \makecell[{{p{\linewidth}}}]{
                \texttt{\tiny{[GM$|$GM]}}
                \texttt{SÖYLE: Ben D nesnesini (6,6) konumuna taşıdım. Sen T nesnesini (3,2) konumuna başarıyla taşıyabildin mi? Kendi tahtandaki güncel durumu ve bir sonraki hamleni paylaşırsan, ben de D nesnesini (6,5) veya (6,4) gibi daha yukarıdaki bir konuma çekerek seninle uyumlu hale getirmeye çalışabilirim.} \\
            }
        }
    }
    & & \\ \\

    \theutterance \stepcounter{utterance}  
    & & \multicolumn{4}{p{0.6\linewidth}}{
        \cellcolor[rgb]{0.9,0.9,0.9}{
            \makecell[{{p{\linewidth}}}]{
                \texttt{\tiny{[P1$\langle$GM]}}
                \texttt{T’ başarıyla (3,2) koordinatlarına taşındı. Güncellenmiş ızgaran şu şekilde:} \\
\texttt{```} \\
\texttt{ 1234567} \\
\texttt{╔════╤══╗} \\
\texttt{║◌◌◌O│M◌║ 1} \\
\texttt{║◌◌TP├──╢ 2} \\
\texttt{║◌S◌◌│◌◌║ 3} \\
\texttt{║◌┌──┼──╢ 4} \\
\texttt{║◌│A◌│D◌║ 5} \\
\texttt{╟─┴┬─┤◌◌║ 6} \\
\texttt{║◌◌│◌│◌◌║ 7} \\
\texttt{╚══╧═╧══╝} \\
\texttt{```} \\
\texttt{Şu anda en fazla 28 turdan 26. turu oynuyorsunuz.} \\
\texttt{Toplamda 16 cezanın 18 kadarını biriktirdiniz.Diğer oyuncu şu mesajı gönderdi:} \\
\\ 
\\ 
\texttt{"Ben D nesnesini (6,6) konumuna taşıdım. Sen T nesnesini (3,2) konumuna başarıyla taşıyabildin mi? Kendi tahtandaki güncel durumu ve bir sonraki hamleni paylaşırsan, ben de D nesnesini (6,5) veya (6,4) gibi daha yukarıdaki bir konuma çekerek seninle uyumlu hale getirmeye çalışabilirim."} \\
\\ 
\\ 
\texttt{Bir sonraki komutun nedir?} \\
            }
        }
    }
    & \\ \\

    \theutterance \stepcounter{utterance}  
    & \multicolumn{4}{p{0.6\linewidth}}{
        \cellcolor[rgb]{0.9,0.9,0.9}{
            \makecell[{{p{\linewidth}}}]{
                \texttt{\tiny{[P1$\rangle$GM]}}
                \texttt{SÖYLE: T nesnesini (3,2) konumuna başarıyla taşıdım. Benim tahtamdaki güncel durum: O (4,1), M (6,1), P (4,2), T (3,2), S (2,3), A (3,5), D (6,5). Nesnelerimiz merkeze doğru oldukça iyi toplandı. Bir sonraki hamlemde A nesnesini (3,4) konumuna çekmeyi deneyeceğim, ancak (3,4) hücresinin boş olduğundan emin olmalıyım. Sen D nesnesini (6,4) konumuna çekebilirsin, böylece koordinatlarımız birbirine daha çok yaklaşır.} \\
            }
        }
    }
    & & \\ \\

    \theutterance \stepcounter{utterance}  
    & & & \multicolumn{2}{p{0.3\linewidth}}{
        \cellcolor[rgb]{0.9,0.9,0.9}{
            \makecell[{{p{\linewidth}}}]{
                \texttt{\tiny{[GM$|$GM]}}
                \texttt{SÖYLE: T nesnesini (3,2) konumuna başarıyla taşıdım. Benim tahtamdaki güncel durum: O (4,1), M (6,1), P (4,2), T (3,2), S (2,3), A (3,5), D (6,5). Nesnelerimiz merkeze doğru oldukça iyi toplandı. Bir sonraki hamlemde A nesnesini (3,4) konumuna çekmeyi deneyeceğim, ancak (3,4) hücresinin boş olduğundan emin olmalıyım. Sen D nesnesini (6,4) konumuna çekebilirsin, böylece koordinatlarımız birbirine daha çok yaklaşır.} \\
            }
        }
    }
    & & \\ \\

    \theutterance \stepcounter{utterance}  
    & & \multicolumn{4}{p{0.6\linewidth}}{
        \cellcolor[rgb]{0.9,0.9,0.9}{
            \makecell[{{p{\linewidth}}}]{
                \texttt{\tiny{[GM$\rangle$P2]}}
                \texttt{Mesajın diğer oyuncuya iletildi.} \\
\\ 
\\ 
\texttt{Şu anda en fazla 28 turdan 27. turu oynuyorsunuz.} \\
\texttt{Toplamda 16 cezanın 18 kadarını biriktirdiniz.Diğer oyuncu şu mesajı gönderdi:} \\
\\ 
\\ 
\texttt{"T nesnesini (3,2) konumuna başarıyla taşıdım. Benim tahtamdaki güncel durum: O (4,1), M (6,1), P (4,2), T (3,2), S (2,3), A (3,5), D (6,5). Nesnelerimiz merkeze doğru oldukça iyi toplandı. Bir sonraki hamlemde A nesnesini (3,4) konumuna çekmeyi deneyeceğim, ancak (3,4) hücresinin boş olduğundan emin olmalıyım. Sen D nesnesini (6,4) konumuna çekebilirsin, böylece koordinatlarımız birbirine daha çok yaklaşır."} \\
\\ 
\\ 
\texttt{Bir sonraki komutun nedir?} \\
            }
        }
    }
    & \\ \\

    \theutterance \stepcounter{utterance}  
    & & & \multicolumn{4}{p{0.6\linewidth}}{
        \cellcolor[rgb]{0.9,0.9,0.9}{
            \makecell[{{p{\linewidth}}}]{
                \texttt{\tiny{[GM$\langle$P2]}}
                \texttt{HAMLE: D, (6,5)} \\
            }
        }
    }
     \\ \\

    \theutterance \stepcounter{utterance}  
    & & & \multicolumn{2}{p{0.3\linewidth}}{
        \cellcolor[rgb]{0.9,0.9,0.9}{
            \makecell[{{p{\linewidth}}}]{
                \texttt{\tiny{[GM$|$GM]}}
                \texttt{HAMLE: D, (6,5)} \\
            }
        }
    }
    & & \\ \\

    \theutterance \stepcounter{utterance}  
    & & & \multicolumn{2}{p{0.3\linewidth}}{
        \cellcolor[rgb]{0.9,0.9,0.9}{
            \makecell[{{p{\linewidth}}}]{
                \texttt{\tiny{[GM$|$GM]}}
                \texttt{D moved to (6, 5)} \\
            }
        }
    }
    & & \\ \\

    \theutterance \stepcounter{utterance}  
    & & \multicolumn{4}{p{0.6\linewidth}}{
        \cellcolor[rgb]{0.9,0.9,0.9}{
            \makecell[{{p{\linewidth}}}]{
                \texttt{\tiny{[P1$\langle$GM]}}
                \texttt{Mesajın diğer oyuncuya iletildi.} \\
\\ 
\\ 
\texttt{Şu anda en fazla 28 turdan 27. turu oynuyorsunuz.} \\
\texttt{Toplamda 16 cezanın 18 kadarını biriktirdiniz.Diğer oyuncu kendi ızgarasında bir nesneyi hareket ettirdi.} \\
\\ 
\\ 
\texttt{Bir sonraki komutun nedir?} \\
            }
        }
    }
    & \\ \\

    \theutterance \stepcounter{utterance}  
    & \multicolumn{4}{p{0.6\linewidth}}{
        \cellcolor[rgb]{0.9,0.9,0.9}{
            \makecell[{{p{\linewidth}}}]{
                \texttt{\tiny{[P1$\rangle$GM]}}
                \texttt{SÖYLE: Diğer oyuncu D nesnesini (6,4) konumuna taşıdı mı? Ben de bir sonraki hamlemde A nesnesini (3,4) konumuna taşımayı deneyeceğim. Lütfen onun hamlesini paylaşır mısın?} \\
            }
        }
    }
    & & \\ \\

    \theutterance \stepcounter{utterance}  
    & & & \multicolumn{2}{p{0.3\linewidth}}{
        \cellcolor[rgb]{0.9,0.9,0.9}{
            \makecell[{{p{\linewidth}}}]{
                \texttt{\tiny{[GM$|$GM]}}
                \texttt{SÖYLE: Diğer oyuncu D nesnesini (6,4) konumuna taşıdı mı? Ben de bir sonraki hamlemde A nesnesini (3,4) konumuna taşımayı deneyeceğim. Lütfen onun hamlesini paylaşır mısın?} \\
            }
        }
    }
    & & \\ \\

    \theutterance \stepcounter{utterance}  
    & & \multicolumn{4}{p{0.6\linewidth}}{
        \cellcolor[rgb]{0.9,0.9,0.9}{
            \makecell[{{p{\linewidth}}}]{
                \texttt{\tiny{[GM$\rangle$P2]}}
                \texttt{D’ başarıyla (6,5) koordinatlarına taşındı. Güncellenmiş ızgaran şu şekilde:} \\
\texttt{```} \\
\texttt{ 1234567} \\
\texttt{╔════╤══╗} \\
\texttt{║◌PMO│◌◌║ 1} \\
\texttt{║◌S◌◌├──╢ 2} \\
\texttt{║A◌◌T│◌◌║ 3} \\
\texttt{║◌┌──┼──╢ 4} \\
\texttt{║◌│◌◌│D◌║ 5} \\
\texttt{╟─┴┬─┤◌◌║ 6} \\
\texttt{║◌◌│◌│◌◌║ 7} \\
\texttt{╚══╧═╧══╝} \\
\texttt{```} \\
\texttt{Şu anda en fazla 28 turdan 28. turu oynuyorsunuz.} \\
\texttt{Toplamda 16 cezanın 18 kadarını biriktirdiniz.Diğer oyuncu şu mesajı gönderdi:} \\
\\ 
\\ 
\texttt{"Diğer oyuncu D nesnesini (6,4) konumuna taşıdı mı? Ben de bir sonraki hamlemde A nesnesini (3,4) konumuna taşımayı deneyeceğim. Lütfen onun hamlesini paylaşır mısın?"} \\
\\ 
\\ 
\texttt{Bir sonraki komutun nedir?} \\
            }
        }
    }
    & \\ \\

    \theutterance \stepcounter{utterance}  
    & & & \multicolumn{4}{p{0.6\linewidth}}{
        \cellcolor[rgb]{0.9,0.9,0.9}{
            \makecell[{{p{\linewidth}}}]{
                \texttt{\tiny{[GM$\langle$P2]}}
                \texttt{SÖYLE: Ben D nesnesini (6,5) konumuna taşıdım. Nesnelerimiz genel olarak merkeze toplandı. Oyunun son turuna geldik. Senin tahtandaki nesnelerin konumları (O:4,1, M:6,1, P:4,2, T:3,2, S:2,3, A:3,4, D:6,5) ile benimkiler (P:2,1, M:3,1, O:4,1, S:2,2, T:4,3, A:1,3, D:6,5) birbirine oldukça yakınlaştı. Eğer sen de nesnelerin konumlarının yeterince yakın olduğunu düşünüyorsan, oyunu bitirebiliriz. SÖYLE: bitti mi?} \\
            }
        }
    }
     \\ \\

    \theutterance \stepcounter{utterance}  
    & & & \multicolumn{2}{p{0.3\linewidth}}{
        \cellcolor[rgb]{0.9,0.9,0.9}{
            \makecell[{{p{\linewidth}}}]{
                \texttt{\tiny{[GM$|$GM]}}
                \texttt{SÖYLE: Ben D nesnesini (6,5) konumuna taşıdım. Nesnelerimiz genel olarak merkeze toplandı. Oyunun son turuna geldik. Senin tahtandaki nesnelerin konumları (O:4,1, M:6,1, P:4,2, T:3,2, S:2,3, A:3,4, D:6,5) ile benimkiler (P:2,1, M:3,1, O:4,1, S:2,2, T:4,3, A:1,3, D:6,5) birbirine oldukça yakınlaştı. Eğer sen de nesnelerin konumlarının yeterince yakın olduğunu düşünüyorsan, oyunu bitirebiliriz. SÖYLE: bitti mi?} \\
            }
        }
    }
    & & \\ \\

    \theutterance \stepcounter{utterance}  
    & & & \multicolumn{2}{p{0.3\linewidth}}{
        \cellcolor[rgb]{0.9,0.9,0.9}{
            \makecell[{{p{\linewidth}}}]{
                \texttt{\tiny{[GM$|$GM]}}
                \texttt{Invalid format: head is not empty: 'SÖYLE: Ben D nesnesini (6,5) konumuna taşıdım. Nesnelerimiz genel olarak merkeze toplandı. Oyunun son turuna geldik. Senin tahtandaki nesnelerin konumları (O:4,1, M:6,1, P:4,2, T:3,2, S:2,3, A:3,4, D:6,5) ile benimkiler (P:2,1, M:3,1, O:4,1, S:2,2, T:4,3, A:1,3, D:6,5) birbirine oldukça yakınlaştı. Eğer sen de nesnelerin konumlarının yeterince yakın olduğunu düşünüyorsan, oyunu bitirebiliriz. '} \\
            }
        }
    }
    & & \\ \\

    \theutterance \stepcounter{utterance}  
    & & & \multicolumn{2}{p{0.3\linewidth}}{
        \cellcolor[rgb]{0.9,0.9,0.9}{
            \makecell[{{p{\linewidth}}}]{
                \texttt{\tiny{[GM$|$GM]}}
                \texttt{Maximum number of penalties exceeded} \\
            }
        }
    }
    & & \\ \\

    \theutterance \stepcounter{utterance}  
    & & & \multicolumn{2}{p{0.3\linewidth}}{
        \cellcolor[rgb]{0.9,0.9,0.9}{
            \makecell[{{p{\linewidth}}}]{
                \texttt{\tiny{[GM$|$GM]}}
                \texttt{Initial states:} \\
\texttt{```} \\
\texttt{Player 1:} \\
\texttt{ 1234567} \\
\texttt{╔════╤══╗} \\
\texttt{║◌◌◌O│PM║ 1} \\
\texttt{║◌◌T◌├──╢ 2} \\
\texttt{║◌◌◌◌│◌◌║ 3} \\
\texttt{║S┌──┼──╢ 4} \\
\texttt{║◌│◌A│◌◌║ 5} \\
\texttt{╟─┴┬─┤◌◌║ 6} \\
\texttt{║◌◌│◌│D◌║ 7} \\
\texttt{╚══╧═╧══╝} \\
\\ 
\\ 
\texttt{Player 2:} \\
\texttt{ 1234567} \\
\texttt{╔════╤══╗} \\
\texttt{║◌◌◌◌│O◌║ 1} \\
\texttt{║P◌M◌├──╢ 2} \\
\texttt{║◌S◌◌│T◌║ 3} \\
\texttt{║◌┌──┼──╢ 4} \\
\texttt{║A│◌◌│◌◌║ 5} \\
\texttt{╟─┴┬─┤◌D║ 6} \\
\texttt{║◌◌│◌│◌◌║ 7} \\
\texttt{╚══╧═╧══╝} \\
\\ 
\texttt{```} \\
            }
        }
    }
    & & \\ \\

    \theutterance \stepcounter{utterance}  
    & & & \multicolumn{2}{p{0.3\linewidth}}{
        \cellcolor[rgb]{0.9,0.9,0.9}{
            \makecell[{{p{\linewidth}}}]{
                \texttt{\tiny{[GM$|$GM]}}
                \texttt{End states:} \\
\texttt{```} \\
\texttt{Player 1:} \\
\texttt{ 1234567} \\
\texttt{╔════╤══╗} \\
\texttt{║◌◌◌O│M◌║ 1} \\
\texttt{║◌◌TP├──╢ 2} \\
\texttt{║◌S◌◌│◌◌║ 3} \\
\texttt{║◌┌──┼──╢ 4} \\
\texttt{║◌│A◌│D◌║ 5} \\
\texttt{╟─┴┬─┤◌◌║ 6} \\
\texttt{║◌◌│◌│◌◌║ 7} \\
\texttt{╚══╧═╧══╝} \\
\\ 
\\ 
\texttt{Player 2:} \\
\texttt{ 1234567} \\
\texttt{╔════╤══╗} \\
\texttt{║◌PMO│◌◌║ 1} \\
\texttt{║◌S◌◌├──╢ 2} \\
\texttt{║A◌◌T│◌◌║ 3} \\
\texttt{║◌┌──┼──╢ 4} \\
\texttt{║◌│◌◌│D◌║ 5} \\
\texttt{╟─┴┬─┤◌◌║ 6} \\
\texttt{║◌◌│◌│◌◌║ 7} \\
\texttt{╚══╧═╧══╝} \\
\\ 
\texttt{```} \\
            }
        }
    }
    & & \\ \\

    \theutterance \stepcounter{utterance}  
    & & & \multicolumn{2}{p{0.3\linewidth}}{
        \cellcolor[rgb]{0.9,0.9,0.9}{
            \makecell[{{p{\linewidth}}}]{
                \texttt{\tiny{[GM$|$GM]}}
                \texttt{* success: False} \\
\texttt{* lose: False} \\
\texttt{* aborted: True} \\
\texttt{{-}{-}{-}{-}{-}{-}{-}} \\
\texttt{* Shifts: 11.00} \\
\texttt{* Max Shifts: 12.00} \\
\texttt{* Min Shifts: 6.00} \\
\texttt{* End Distance Sum: 10.48} \\
\texttt{* Init Distance Sum: 20.21} \\
\texttt{* Expected Distance Sum: 29.33} \\
\texttt{* Penalties: 19.00} \\
\texttt{* Max Penalties: 16.00} \\
\texttt{* Rounds: 27.00} \\
\texttt{* Max Rounds: 28.00} \\
\texttt{* Object Count: 7.00} \\
            }
        }
    }
    & & \\ \\

\end{supertabular}
}

\end{document}
